% concatenated from Short_Main_0826.tex
\documentclass{article}
\usepackage[nonatbib,preprint]{neurips_2026}
% \usepackage[nonatbib,sglblindworkshop]{neurips_2026}

\usepackage[utf8]{inputenc}
\usepackage[T1]{fontenc}
\usepackage{amsmath,amssymb,amsthm,mathtools}
\usepackage{enumitem}
\usepackage{microtype}
\usepackage{xcolor}
\usepackage{hyperref}
\usepackage{graphicx}
\usepackage{subcaption}



\usepackage[ruled,vlined,linesnumbered]{algorithm2e}
\SetKwComment{Comment}{$\triangleright$\ }{}
\usepackage{dsfont}

\usepackage{natbib}
\usepackage{aliascnt}

\newtheorem{theorem}{Theorem}[section]

\newaliascnt{lemma}{theorem}
\newtheorem{lemma}[lemma]{Lemma}
\aliascntresetthe{lemma}

\newaliascnt{proposition}{theorem}
\newtheorem{proposition}[proposition]{Proposition}
\aliascntresetthe{proposition}

\newaliascnt{corollary}{theorem}
\newtheorem{corollary}[corollary]{Corollary}
\aliascntresetthe{corollary}

\newaliascnt{assumption}{theorem}
\newtheorem{assumption}[assumption]{Assumption}
\aliascntresetthe{assumption}

\theoremstyle{definition}

\newaliascnt{definition}{theorem}
\newtheorem{definition}[definition]{Definition}
\aliascntresetthe{definition}

\newaliascnt{remark}{theorem}
\newtheorem{remark}[remark]{Remark}
\aliascntresetthe{remark}

\makeatletter
\newtheorem*{rep@theorem}{\rep@title}
\newcommand{\newreptheorem}[2]{%
\newenvironment{rep#1}[1]{%
\def\rep@title{#2 \ref{##1}}%
\begin{rep@theorem}}%
{\end{rep@theorem}}}
\makeatother

\newreptheorem{theorem}{Theorem}
\newreptheorem{lemma}{Lemma}
\newreptheorem{corollary}{Corollary}

\usepackage[capitalize]{cleveref}

\crefname{theorem}{Theorem}{Theorems}
\Crefname{theorem}{Theorem}{Theorems}

\crefname{lemma}{Lemma}{Lemmas}
\Crefname{lemma}{Lemma}{Lemmas}

\crefname{proposition}{Proposition}{Propositions}
\Crefname{proposition}{Proposition}{Propositions}

\crefname{corollary}{Corollary}{Corollaries}
\Crefname{corollary}{Corollary}{Corollaries}

\crefname{assumption}{Assumption}{Assumptions}
\Crefname{assumption}{Assumption}{Assumptions}

\crefname{definition}{Definition}{Definitions}
\Crefname{definition}{Definition}{Definitions}

\crefname{remark}{Remark}{Remarks}
\Crefname{remark}{Remark}{Remarks}

\newcommand{\E}{\mathbb E}
\newcommand{\N}{\mathbb N}
\newcommand{\pro}{\mathbb P}
\newcommand{\conv}{\mathrm{conv}}
\newcommand{\tagc}{\mathrm{tag}}
\newcommand{\init}{\mathrm{init}}
\newcommand{\newb}{\mathrm{new}}
\newcommand{\LRU}{\mathrm{LRU}}
\newcommand{\one}{\mathds{1}}

\newcommand{\Exp}{\text{Exp}}
\newcommand{\Pois}{\text{Poisson}}
\newcommand{\Erlang}{\text{Erlang}}
\newcommand{\SLRU}{\mathrm{SLRU}}
\newcommand{\Op}{O_{\mathbb P}}
\newcommand{\yhy}[1]{[\textcolor{blue}{Heyuan: #1}]}
\newcommand{\chutong}[1]{[\textcolor{red}{Chutong: #1}]}


\title{Multi-Turn LLM Conversations under the Least-Recently-Used Policy: Mean-Field Asymptotics and Hit Ratio Approximation}

\author{
  Heyuan Yao \\
  Northwestern University\\
  Evanston, IL 60208 \\
  \texttt{heyuanyao@u.northwestern.edu}
  \And
  Chutong Gao\textsuperscript{*} \\
  Purdue University \\
  West Lafayette, IN 47907\\
  \texttt{gao945@purdue.edu}
  \AND
  Yuan Lyu\textsuperscript{*} \\
  HKUST\\
  Hong Kong SAR\\
  \texttt{ylyuad@connect.ust.hk}
  \And
  Izzy Grosof \\
  Northwestern University\\
  Evanston, IL 60208 \\
  \texttt{izzy.grosof@northwestern.edu}
  \And
  David Simchi-Levi \\
  Purdue University \\
  West Lafayette, IN 47907\\
  \texttt{dslevi@purdue.edu}
  \AND
}

\date{}
\begin{document}
\maketitle

\vspace{-5.5em}
\begin{center}
\small \textsuperscript{*}Equal contribution.
\end{center}
\vspace{0.5em}

\begin{abstract}
The major workloads in modern large language model (LLM) serving systems have shifted from single-shot LLM calls to multi-turn conversations, where new responses are generated based on the whole conversation history across all previous turns.    
The hit ratio, i.e., the average fraction of KV caches accessed directly from existing caches stored in high-bandwidth memory (HBM), is hence a crucial metric that governs system performance. Estimating the hit ratio is a highly nontrivial task due to the complex system dynamics, where the KV cache prefixes grow with turns and some must be evicted due to finite memory capacity. We formulate the system as a multi-turn conversation model under the least-recently-used (LRU) policy. Through a mean-field asymptotic framework, we prove that as the conversation arrival rate and the memory capacity grow proportionally to infinity, the hit ratio converges to a closed-form limit. Based on the characterization of the limit, we further propose a practical hit ratio estimator, and validate its accuracy by real LLM serving experiments on the Qwen3-8B model implemented on Ascend NPUs. Our results provide a theoretical foundation for the analysis of multi-turn LLM serving systems and a practical guideline for memory capacity provisioning.
\end{abstract}

\vspace{-0.5em}

\vspace{-0.5em}
\vspace{-0.5em}

\section{Introduction}
\label{sec: Introduction}
\vspace{-0.5em}
\vspace{-0.2em}




Multi-turn interactions have become a core workload form in modern large language model (LLM) services, following the prevalence of coding agents as the most commercially validated LLM application (\cite{github2025codingagent, anthropic2026claudecode}).
%GitHub reports that its coding agent helped developers merge more than one million pull requests within five months of launch (\cite{github2025codingagent}), while Anthropic reports that Claude Code has surpassed \$2.5 billion in run-rate revenue (\cite{anthropic2026claudecode}). 
Within one \textit{conversation}, 
coding agents must carry all conversation history (including previous system instructions, history prompt-response pairs, and tool-calling feedback) for new response generation, and such \textit{multi-turn} LLM calls create prolonged and growing conversation history. A key strategy to hedge against the correspondingly growing compute is context caching, that is, storing and reusing the already generated key-value (KV) caches of the conversation history in the memory or disk, such that new prompts sharing the same history can avoid repeated prefill and thus massively reduce time-to-first-token (TTFT) and inference costs. For example, DeepSeek reports that context caching reduces the TTFT of repetitive 128K-token inputs from 13s to 500ms, and consequently, it charges a prompt that successfully hits the cached content $1/10$ of the original price (\cite{deepseek2024contextcache}). 

To leverage the maximum value of context caching, it is important to design an LLM serving system that
maximizes the long-run average \textit{hit ratio}, defined as the fraction of KV caches directly reused from the HBM among all KV caches in a generic conversation. A prerequisite for the optimization problem is the following performance analysis problem, which is the focus of this paper:

\vspace{-0.5em}

\begin{quote}
\textit{How to characterize (or approximate) the hit ratio as a function of given system primitives, including the conversation arrival rate, the interarrival-time distribution of prompts within one conversation, the workload distributions, and the offered HBM capacity?}
\end{quote}
\vspace{-0.5em}
\vspace{-0.5em}

\paragraph{The least-recently-used (LRU) eviction policy.}
A conversation history can be reused only if its KV caches remain resident in the HBM memory when a new prompt within the same conversation arrives. However, the HBM memory of the processing unit has only finite capacity and thus inevitably forces eviction of its resident conversations. Throughout the paper we assume the system operates under the LRU eviction policy widely adopted in multi-turn KV-cache management due to its simplicity and empirical success (\cite{wang2025kvcache, zhang2026tail}). Under LRU, when the amount of KV caches exceeds the HBM capacity, conversations (with its all history KV caches) are evicted sequentially from the least-recently-used one to the most-recently-used one until the capacity constraint is restored.
%a reuse hit is governed by a race between the conversation's return time and its cache eviction time: the former depends on inter-turn times and conversation lifetimes, whereas the latter depends on conversation arrivals, prefix growth, and cache capacity. Under-provisioning evicts reusable states prematurely, while over-provisioning consumes expensive memory for diminishing hit ratio gains. The hit ratio is thus jointly determined by workload dynamics, serving load, and provisioned capacity.

%This interaction creates an analytical provisioning problem. Existing systems provide block-based KV management, hierarchical storage, and KV offloading, while trace-driven simulators can evaluate candidate configurations (\cite{gao2024cachedattention,qin2025mooncake,agrawal2024vidur}). Yet cache sizing still relies largely on trace replay, simulation, and empirical search. What is missing is a direct mapping from observable multi-turn workload statistics to the LRU hit-ratio curve, and from a target hit ratio back to the required capacity. We therefore ask:
%\begin{quote}
%    \emph{Given a stochastic multi-turn LLM workload and a model-specific KV-memory footprint, what is the cache capacity required under LRU to achieve a target  hit ratio?}
%\end{quote}
\vspace{-0.5em}

\paragraph{Overview of our results and contributions.}
We make three main contributions.
\vspace{-0.5em}

\begin{itemize}
    \item \textbf{Modeling framework.} We propose a parsimonious multi-turn conversation model (MCM) for multi-turn LLM serving under LRU eviction policy, which captures key components that govern the stochastic system dynamics.
    \item \textbf{Mean-field asymptotics.} We prove the convergence of the hit ratio to an  explicit limit as the conversation arrival rate and the HBM capacity grows proportionally to infinity.
    \item \textbf{Estimator and validation.} We derive a theoretically-motivated hit ratio estimator and validate its accuracy in Qwen3-8B model serving experiments implemented on Ascend 910B2 NPUs fed by conversations from the ShareGPT dataset.
\end{itemize}
% \paragraph{Contribution.}
%To answer this question, we generalize the stochastic conversation model of~\citet{zhang2026tail} to a \emph{Multi-Turn Conversation Model} (MCM) and develop a block-weighted mean-field analysis of LRU. From the conversation arrival rate, inter-turn distribution, conversation length, and per-turn KV growth, we characterize the tagged-conversation displacement process and derive the characteristic time and  hit ratio. This mapping can be inverted to compute the capacity required for a target hit ratio. These statistics can be estimated from serving traces, while the model's KV-memory footprint translates the predicted number of blocks into physical memory. We further develop a theoretically-motivated estimator and validate it on real-world multi-turn conversation traces. Although our experiments use an HBM-resident cache, the analysis depends on workload statistics, effective cache capacity, and eviction dynamics rather than a specific storage medium, and may serve as a first-order approximation for hierarchical and offloaded KV-cache systems.

\vspace{-0.2em}

\vspace{-0.5em}
\vspace{-0.5em}

\section{The Model}
\label{sec: Model Settings-short}
\vspace{-0.5em}
\vspace{-0.2em}

In this section, we formulate the MCM under the LRU policy and the corresponding  hit ratio. Additional model details, formal policy definitions, and illustrations are provided in \cref{sec: Model Settings}. 

% \paragraph{Notation.} We denote by $\N$ the set of natural numbers, and write $\N_+ := \N \backslash \{0\}$. 
\vspace{-0.5em}

\paragraph{The multi-turn conversation model.}
We first introduce the MCM in detail.
In LLM practice, KV caches are managed as blocks in the HBM memory, each having a fixed capacity measured in token positions\footnote{For instance, in our vLLM-Ascend experiments, each full KV block represents 128 token positions.}. We consider a system where the HBM can store at most $N$ KV blocks, referred to as the \textit{HBM capacity}. We assume the system is implemented with a Prefill-Decode disaggregation structure with one prefill instance, where the prefill and decode workloads are processed by different machines (\textit{prefiller} and \textit{decoder}). The HBM capacity $N$ under our consideration is that of one prefiller. New \textit{conversations} arrive according to a Poisson process with rate $\lambda_0$.  
We assume that the numbers of turns $M$ of different conversations are i.i.d.
Upon arrival of a conversation with its initial prompt tokens, turn-$1$ starts and the prefiller generates $X_1$ KV blocks for the prompt. Its corresponding response is immediately generated by one decoder, whose tokens, but not the KV caches\footnote{The current Mooncake infrastructure (\cite{qin2025mooncake}) supports the KV block transfer only from a prefiller to a decoder, but not in the reverse direction.}, are sent back to the prefiller. This marks the end of turn-$1$. After a random interarrival time $W$ (i.i.d. across all turns and all conversations), a new prompt of the same conversation arrives, and a new turn starts.  
At the start of turn-$j$, $j=2, \ldots, M$, the prefiller immediately generates $Y_{j-1}$ KV blocks for the response to the last prompt from turn-$(j-1)$, and $X_j$ KV blocks for the new prompt of turn-$j$. Then the corresponding KV blocks are transferred to the decoder, to generate the turn-$j$ response. Again, the tokens of this response, but not the KV blocks, are sent back to the prefiller. This marks the end of turn-$j$. Therefore, the KV block increment of turn-$j$ $Z_j$ is defined by $Z_1:=X_1 ~\text{ and }~ Z_j:=Y_{j-1} + X_j, j\geq2$,
where we assume $Z_j\in\N_+$. The history block size immediately after the turn-$j$ prompt is therefore defined as $R_j:=\sum_{\ell=1}^{j}Z_\ell$. \cref{fig: KV content increments example} illustrates the growth of the KV content. In \cref{fig: Multi-turn examples}, we use a running example to better illustrate the model setting.


Our analytical results for the MCM are established under the following additional assumptions. 
First, the turn-wise block increments $\{Z_{j}\}_{j\geq 1}$, the interarrival times $W$ between each turn, and the turn number $M$ are mutually independent, while dependence across block increments of different turns is allowed. In addition, the cumulative number of KV blocks generated at the last turn of a conversation, defined as
\begin{equation} \label{eq:B}
   B:=R_M=\sum_{\ell=1}^{M} Z_\ell, 
\end{equation}
satisfies $\E [\exp( \theta_0 B)]<\infty$ for some $\theta_0>0$. Finally, the interarrival time $W$ of prompts within one conversation has a Lipschitz continuous cumulative distribution function (CDF) $F$.

\vspace{-0.5em}

\paragraph{The least-recently-used policy and the  hit ratio.}
Consider a generic conversation with $M$ turns and history content sizes $\{R_j\}_{j=1}^M$. If the conversation continues from turn $j$ to turn $j+1$, then the $R_j$ history blocks become reuse-eligible at the next prompt, as illustrated in \cref{fig: LLM multi turn example}. If any prefix block cannot be recovered when its next-turn prompt arrives, all suffix blocks of this prefix block cannot be retrieved from the HBM\footnote{In reality, these suffix blocks must be retrieved from lower-level storage or reprefilled, which causes more computation and transfer cost than the local HBM reuse. }. 

The MCM under our consideration operates under the LRU policy. The system maintains an ordered resident conversation list $\mathcal L$, from the most-recently-used (MRU) end to the LRU end. The list is updated whenever a new prompt is processed and the corresponding conversation content is inserted or refreshed in HBM. Suppose an arriving prompt belongs to a resident conversation $\theta\in\mathcal L$. The conversation label $\theta$ is moved to the MRU end, while its previously stored content history is replaced by the enlarged content containing the newly generated blocks. If the prompt belongs to a nonresident conversation $\theta\notin\mathcal L$, the conversation is inserted directly at the MRU end, together with its available history content and the newly generated blocks.

The HBM can store at most $N$ KV blocks. Whenever inserting or updating a conversation content causes the total resident KV blocks to exceed this capacity, the system evicts complete conversation contents sequentially from the LRU end until the capacity constraint is restored\footnote{In practical implementations, LRU eviction may operate at the block level and therefore allow partial-prefix eviction. We discuss these implementation differences in the appendix.}. The whole-content LRU policy under our consideration is described by Algorithm \ref{alg: Whole-Content LRU} in the Appendix, and the hit ratio under LRU admits the following definition. 
    For each turn $j\geq 2$, define the LRU hit indicator
    \[H_{j} := \one \left\{ \text{the history content of size $R_{j-1}$ is resident in HBM when the turn-$j$ prompt arrives} \right\}.
    \]
    The cache hit ratio is defined as  $h  :=\E [\sum_{j=2}^{M}R_{j-1}H_{j}]/{A_R}$, where $A_R : =\E [\sum_{j=1}^{M} R_j]$ represents the expected cumulative history content workload over the lifetime of a conversation, and the numerator represents the portion of this workload that is reused through direct HBM lookup.


\vspace{-0.5em}
\vspace{-0.5em}

\paragraph{The mean-field asymptotic regime.}


An explicit characterization of the hit ratio $h$ is rendered intractable due to the complex system dynamics, so we resort to a mean-field scaling regime where the conversation arrival rate scales proportionally with the HBM capacity $N$, and we characterize the mean-field limit of the hit ratios.
Specifically, the system indexed by $N$ has HBM capacity $N$ and conversation arrival rate $\lambda_0^{(N)}:=N\lambda_0$. The other system primitives, i.e., the prompt interarrival-time distribution $F$, the turn-number distribution $M$, and the block-increment sequence $\{(X_j,Y_j)\}_{j\geq1}$, are kept unchanged. Under this scaling, both the HBM capacity and the aggregate workload arrival rate grow linearly with $N$, and we will show in the following section that this scaling yields a nondegenerate mean-field limit for the hit ratio.
\vspace{-0.5em}
\vspace{-0.5em}

\section{Main results}
\label{sec: main results-short}
\vspace{-0.5em}
\vspace{-0.2em}

In this section, we present our main theoretical result. We establish the convergence of the hit ratios for a sequence of MCMs operating under LRU and indexed by their HBM capacities $N = 1, 2, \ldots$. Based on the mean-field convergence results, we develop a theoretically-motivated estimator for the hit ratio, and we validate its accuracy through real LLM serving experiments. 
%Unless otherwise stated, all results assume a stationary system.
For more details, we refer our readers to \cref{sec: main results}. All proofs in this paper appear in \cref{sec: Proofs}.

We consider the normalized displacement process of a tagged conversation, i.e., the cumulative KV block mass (normalized by N ) inserted or refreshed ahead of the tagged content in the LRU order since its last access. To state the main result,  we define the mean displacement curve $d:\mathbb R_+  \rightarrow \mathbb R_+$, which is shown to be the mean of the normalized displacement process (see \cref{lem: Mean Displacement Curve}), as
\begin{equation}
d(t):=\lambda_0\E[B]t+\lambda_0A_H\int_0^t\bar F(u)\,du,
\end{equation}
where $B$ is given in \eqref{eq:B}, $A_H:=\E[\sum_{j=1}^{M-1} R_j]$ denotes the tight upper bound of the expected amount of reusable history content workload, and $\bar{F}$ is the tail probability of the prompt interarrival time $W$.
Note that $d$ is strictly increasing, and has a well-defined inverse $d^{-1}$. We then define the mean-field characteristic time $T_C:=d^{-1}(1)$, i.e., the time at which the mass fills the entire cache.

We are now ready to state the main result. For the system indexed by $N$, we let $h^{(N)}$ denote the  hit ratio and define the \textit{mean-field hit ratio} as $h ^{(\infty)} :=(A_H/{A_R})F(T_C) = (\E[\sum_{j=1}^{M-1} R_j] / \E [\sum_{j=1}^{M} R_j]) F(T_C)$. %\chutong{$A_H$ undefined}

\begin{theorem}[Mean-Field Limit of the Hit Ratio]
    \label{thm: Hit Ratio Convergence in the MCM-LRU model- Short-short}
    Under the MCM assumptions, we have $h ^{(N)}\rightarrow h ^{(\infty)}$ as $N\rightarrow \infty.$
\end{theorem}
\vspace{-0.5em}
\vspace{-0.3em}

\paragraph{theoretically-motivated hit ratio estimator.}
In real multi-turn LLM conversation systems, the last KV block generated at each turn is usually partially filled, which cannot be reused. We then modify the MCM to a practical version, where we further assume that the last block of each turn is partially filled and hence unhashable for prefix matching. Excluding one such block per non-final turn gives the expected reusable workload $A_H^{\mathrm{hash}}:=A_H-\E[M-1].$ Applying the same mean-field argument to the practical MCM, we have that the hit ratio $h_{ \mathrm{prac}}^{(N)}$ converges to the mean-field hit ratio $h_{ \mathrm{prac}}^{(\infty)}:={(A_H^{\mathrm{hash}}}/{A_R})\,F(T_C^{\mathrm{prac}})$, where $T_C^{\mathrm{prac}}$ is the practical characteristic time, which is the  solution to the  equation $d_{\mathrm{prac}}(t):=\lambda_0 \E[B]t+\lambda_0 A_H^{\mathrm{hash}}\int_0^t \bar F(u) du =1.$

Based on our convergence results, we propose a theoretically-motivated mean-field estimator for the hit ratio. Suppose that we have the estimates $\hat \lambda_0$, $\hat A_R, \widehat A_H^{\mathrm{hash}}, \widehat{\E[B]}, \widehat F$ of $\lambda_0$, $A_R, A_H^{\mathrm{hash}}, {\E[B]}, F$, respectively, and we have HBM capacity $N$. Then our estimator is given by 
\begin{equation} \label{eq:estimator}
    \widehat h:={\widehat A_H^{\mathrm{hash}}}/{\widehat A_R} \, \widehat F(\widehat T_C^{\mathrm{prac}}),
\end{equation}
where $\widehat T_C^{\mathrm{prac}}$ is the solution to $\hat \lambda_0 \widehat{\E[B]}\widehat T_C^{\mathrm{prac}}+\hat \lambda_0 \widehat A_H^{\mathrm{hash}}\int_0^{\widehat T_C^{\mathrm{prac}}}\widehat{\bar F}(u)\,du=N.$ Here $\hat \lambda_0$ denotes the unscaled conversation arrival rate, so the threshold
is $N$ rather than $1$.


\vspace{-0.5em}
\vspace{-0.5em}

\section{LLM Serving Experiments}
\label{sec: LLM Serving Experiments}
% \chutong{model - card - dataset. Cut to 1/3 of the current length. Cite ShareGPT. Put our code online and refer to it. Second paragraph: report results. }
\vspace{-0.5em}
\vspace{-0.2em}


In this section, we validate our estimator \eqref{eq:estimator} using public multi-turn ShareGPT conversations (\cite{sharegpt_vicuna_unfiltered}). Because each trace contains both prompts and responses, we use Prefill-Decode disaggregation to avoid inconsistencies between the generated responses and subsequent recorded prompts. We deploy Qwen3-8B on 5 Ascend 910B2 NPUs, with one NPU as the prefiller.
% \chutong{No need to appear here; moved to appendix} For each prompt, the decoder generates the same number of tokens as the recorded response, but its output is discarded. Instead, we append the recorded response to the next-turn prompt to send to the prefiller after a sampled interarrival time. 
We assign 40 GB of prefiller HBM for the KV cache, such that the block capacity is $N=2,275$. We use 5,000 conversations for measurement under Poisson conversation arrivals. More details on the experimental setting are deferred to \cref{sec: LLM-Inference Experiment and Empirical Results}.

In Experiment 1, we fix $\lambda_0=1.5$ and vary the target mean prompt interarrival time $\mu_F$, using an exponential clock. However, a new prompt is submitted when both the sampled exponential clock rings and the previous decoding finishes. The prior estimator uses $F(t)=1-exp(-\frac{t}{\mu_F})$, while the posterior estimator replaces $\mu_F$ by the measured mean interarrival time $\hat{\mu}_F$. We report the absolute error $\mathrm{AE}=|\text{estimated hit ratio}-\text{actual hit ratio}|$ and relative error $\mathrm{RE}=\mathrm{AE}/\text{actual hit ratio}\times100\%$. In Experiment 2, we fix $\mu_F=135$s, under which $\hat{\mu}_F$ is closest to the target, and vary the conversation arrival rate from $0.5$ to $2.0$. Results are shown in \cref{fig: E1-1 hit ratio results,fig: E1-1 Error Analysis,fig: E1-2 hit ratio results,fig: E1-2 Error Analysis}. Both estimators capture the decrease in empirical hit ratio when either $\mu_F$ or the $\lambda_0$ increases. In addition, the AE remains below $0.02$, and the RE is below $10\%$ when the realized interarrival time closely matches its target.


\vspace{-0.5em}
\vspace{-0.5em}

\section{Conclusion}
\vspace{-0.5em}
\vspace{-0.2em}

We studied the characterization of the hit ratio in LLM serving systems with multi-turn conversations.   
We proposed a parsimonious multi-turn conversation model under the widely adopted least-recently-used policy, which captures several key components governing the real system dynamics. Through a mean-field asymptotic framework, we proved the convergence of the hit ratios to a certain mean-field hit ratio limit as the HBM capacity $N\to \infty$. Based on our characterization of the mean-field limit, we designed a hit ratio estimator for practical LLM serving systems, and validated its accuracy via LLM serving experiments with the Qwen3-8B model implemented on Ascend 910B2 NPUs. Our work provides a theoretical foundation for analyzing multi-turn LLM services, and provides practical system-design guidelines on memory provisioning that achieves a target hit ratio. 


% \section{Conclusion}




% \chutong{Some standarization for terminology:  1. We write ``least-recently-used policy'' with ``-''; abbreviation is introduced where they were mentioned at the first time and we use the abbreviation thereafter, e.g., multi-turn conversation model (MCM).... and we use MCM only after that. 2. Call $X_i$ the number of prompt blocks and $Y_i$ the number of response blocks. (these are more easy to understand; avoid prefill block and decode block) }


\bibliographystyle{plainnat}
\bibliography{reference}

\appendix

\section{Prior Work}
\label{sec: prior work}


\subsection{Caching Systems and KV-Cache Management for Multi-turn LLM Serving}

% Recent LLM serving systems support cross-request KV-cache reuse at several layers. CachedAttention studies KV reuse across turns, Mooncake expands effective cache capacity through a tiered and disaggregated architecture, and production-workload characterization shows that KV-cache reuse intervals and probabilities vary substantially across applications (\cite{gao2024cachedattention,qin2025mooncake,wang2025kvcache}). These mechanisms have also entered practical inference platforms: TensorRT-LLM manages limited KV-cache blocks using LRU and priority-based eviction and supports KV-aware routing, while the DeepSeek API reports and prices cache-hit and cache-miss prompt tokens separately (\cite{nvidia2025kvcachereuse,deepseek2024contextcache}). Because different requests may reuse substantially different numbers of KV blocks, a request-level hit indicator does not capture the volume of reuse. We therefore adopt a block-weighted hit ratio, which measures the fraction of reusable KV-state workload served from HBM.



Recent caching systems reveal several workload characteristics that are particularly relevant to multi-turn LLM serving. In conventional caching systems, some caches' highly heterogeneous object sizes have already been observed in production workloads and incorporated into cache admission and resource allocation (\cite{berger2017adaptsize,berger2018robinhood,berg2020cachelib}). In LLM serving, KV states are typically managed at the block level (\cite{kwon2023pagedattention}), while block reuse may span multiple conversation turns and even multiple storage tiers (\cite{deepseek2024contextcache,gao2024cachedattention,qin2025mooncake}). Production traces further show substantial heterogeneity in KV-cache reuse intervals, reuse probabilities, and reusable prefix lengths across requests \cite{wang2025kvcache}. Practical serving platforms therefore employ mechanisms such as LRU- or priority-based eviction, prefix matching, and KV-aware routing under limited cache capacity (\cite{nvidia2025kvcachereuse}). Moreover, a distinctive feature of multi-turn LLM workloads is that the cacheable blocks associated with the same conversation grow across successive accesses, rather than remaining a fixed-size object.

\subsection{Previous Analyses of LRU and Hit Ratio}
\label{ssec: Previous Analyses of LRU and Hit Ratio}
    
We now review prior analytical work on the least-recently-used (LRU) cache management policy. Existing analyses have mainly focused on web and content caching systems, where a fixed population of cacheable objects is repeatedly requested according to independent, stationary, or temporally correlated request processes. A common objective is to characterize the miss or hit probability of LRU in large-cache regimes, often through asymptotic approximations that replace the complicated LRU ordering dynamics with a more tractable working-set or time-to-live (TTL) representation. One of the earliest asymptotic analyses of LRU can be found in \cite{fagin1977asymptotic}, which showed that the expected LRU miss ratio can be asymptotically approximated by the miss ratio of a corresponding working-set model. This working-set model characterizes cache residency through a fixed window of recent page references. A closely related line of work analyzes LRU through a deterministic characteristic time that approximates the random eviction time of a cached object. \citet{che2002hierarchical} introduced this characteristic time approximation for hierarchical web caches, which was later widely referred to as the Che approximation. A mathematical explanation for its accuracy was provided in \cite{fricker2012versatile} by relating the eviction of a tagged object to the cumulative number of distinct objects requested after its last access. In \cite{martina2014unified}, this characteristic time decoupling principle was further extended to general renewal traffic, multiple cache management policies, and interconnected caches. Subsequent work established stronger asymptotic foundations for characteristic time and TTL approximations under more general request processes, including shot-noise request models with temporal locality (\cite{leonardi2017modeling}) and independent stationary and ergodic request processes (\cite{jiang2018convergence}). In addition, \citet{gast2017ttl} connected TTL approximations for LRU variants to deterministic transient mean-field ODE limits under the IRM, whose fixed points recover the proposed TTL approximations.

Our analysis builds on a similar general principle. We represent LRU eviction through the cumulative workload that moves ahead of a tagged conversation content and identify a deterministic characteristic time in the mean-field limit. However, multi-turn LLM conversations introduce several challenges that are absent from previous theoretical work, which focuses on traditional object-based caching models. First, cacheable KV blocks have finite reuse lifetimes. A block can only be reused by subsequent turns of its corresponding conversation. Once the conversation terminates, its blocks may still occupy HBM but can never generate another cache hit. Moreover, the server does not know in advance when a conversation will terminate, which contrasts with classical IRM settings with a fixed population of objects and time-invariant request probabilities. Second, the reusable history associated with a conversation grows over time, so a later access can touch a nondecreasing number of KV blocks. As a result, the displacement process must track the amount of KV block workload moved ahead of the tagged content, while the size of the tagged content itself is also random and increases with the turn state. Finally, request-level hit probability is insufficient for quantifying cache effectiveness in LLM serving because different prompts may reuse substantially different numbers of historical KV blocks. Inspired by the byte hit ratio used in content delivery network (\cite{berger2017adaptsize}), we therefore define a hit ratio that weights successful reuse by the number of KV blocks served from HBM, which further quantifies the computational and data-movement savings benefited by LRU caching.



\section{Model details and formal definitions}
\label{sec: Model Settings}

In this section, we formulate the multi-turn conversation model (MCM) and introduce the whole-content Least-Recently Used (LRU) policy. We begin with basic notation in \cref{ssec: Basic Notations}. In \cref{ssec: Generalized SCM}, we define the MCM for multi-turn LLM inference. We then introduce the  hit ratio and characterize its intrinsic upper bound in \cref{ssec: The  Hit Ratio}. In \cref{ssec: The LRU Policy}, we specify the whole-content LRU policy used in our theoretical analysis. The mean-field scaling is introduced in \cref{ssec: Our Mean-Field Model}. Finally, in \cref{ssec: tagged-conversation Displacement Process and Characteristic Time}, we define the tagged-conversation displacement process and its characteristic time. These two terms will be used as the key stochastic objects in our mean-field analysis in \cref{sec: main results}. Their limiting behavior is important for our subsequent hit ratio convergence analysis.

\subsection{Basic Notations}
\label{ssec: Basic Notations}

We let $\mathbb N$ denote the set of nonnegative integers and $\mathbb N_+$ denote the set of positive integers. For a cumulative distribution function $F$, its survival function is defined as $\bar F$ with $\bar F(t):=1-F(t)$. 

\subsection{Multi-Turn Conversation Model}
\label{ssec: Generalized SCM}

We first introduce our multi-turn conversation model (MCM). In LLM practice, KVs are managed as blocks in the system, the number of which can be modeled to follow some integer distribution. For instance, in our experiment, 128 tokens are arranged in a block, which always has a constant memory requirement. We consider a system with an  HBM  KV block memory capacity $N\in\mathbb N$, i.e., the HBM can store at most $N$ KV blocks. New conversations arrive according to a Poisson process of rate $\lambda_0$. Each conversation enters the system and immediately produces a prompt corresponding to $X_1$ KV blocks, and the LLM generates a sentence corresponding to $Y_1$ KV blocks. We assume that the number of turns of a conversation follows some positive integer distribution $M$. At turn-$j$, the numbers of prompt blocks and response blocks generated follow a joint distribution $(X_j, Y_j) \in \N^2$. 


Our model is motivated by the Prefill-Decode disaggregation setting, where the KV block composition at successive turns is visualized in \cref{fig: KV content increments example}. We define the turn-wise content increment as $Z_1:=X_1$ and
\begin{equation*}
    Z_j:=X_j+Y_{j-1},\qquad j\geq2.
\end{equation*}
We assume $Z_j\in\N_+$, so that every turn increases the history content by at least one KV block. The history content size immediately after the turn-$j$ prompt is therefore defined as
\begin{equation}
    \label{equL The history content size R_j}
    R_j:=\sum_{\ell=1}^{j}Z_\ell=\sum_{\ell=1}^{j}X_\ell+\sum_{\ell=1}^{j-1}Y_\ell.
\end{equation}
Thus, $R_j$ contains all prompt KV blocks up to turn $j$ and all response KV blocks generated up to turn $j-1$. This entire history content is required to process the turn-$j$ prompt and generate the corresponding response.   

\paragraph{A running example.} \cref{fig: LLM multi turn example} illustrates a running example of how KV blocks in one turn are managed. Lowercase letters denote tokens, and uppercase letters denote block numbers. The tokens at the start of turn-$3$ are $(x_1, y_1, x_2, y_2, x_3)$. All history KV blocks (or ``prefix KV blocks'') $(X_1,Y_1, X_2)$ are retrieved and reused from HBM, based on which the prefiller generates $Y_2$ KV blocks for the response $y_2$ from the last turn, and $X_3$ blocks for the new prompt $x_3$. All KV blocks $(X_1, Y_1, X_2, Y_2, X_3)$ are passed to a decoder to generate the new response $y_3$ for the prompt $x_3$. But only the tokens $y_3$, not the KV blocks $Y_3$, are sent back to the prefiller for turn-$4$'s compute.

\begin{figure}[ht] 
    \centering
    \begin{subfigure}[t]{0.39\linewidth}
        \centering
        \includegraphics[width=\linewidth]{KV_content_increments_example.png}
        \caption{Multi-turn prefix cache reuse example}
        \label{fig: KV content increments example}
    \end{subfigure}
    \hfill
    \begin{subfigure}[t]{0.59\linewidth}
        \centering
        \includegraphics[width=\linewidth]{LLM_multi_turn_example.png}
        \caption{Turn-Level KV cache workflow example}
        \label{fig: LLM multi turn example}
    \end{subfigure}
    \caption{(a) Illustration of the KV block composition at successive turns, under the Prefill-Decode disaggregation setting. (b) An example of how the KV blocks in one turn are managed, assuming the full history KV blocks reside in HBM.}
    \label{fig: Multi-turn examples}
\end{figure}



\begin{remark}
    The treatment of response KV blocks may differ across several state-of-the-art LLM serving architectures. The HBM KV block increment dynamics described above are most relevant in the Prefill-Decode disaggregation architecture in which response KV blocks are released after decoding. In this case, when the conversation receives the next-turn prompt, the corresponding response history must be reconstructed on the prefill side.

    Under some Prefill-Decode disaggregation frameworks, however, response KV blocks may be transferred back to the prefill-side worker through a ``write-back" mechanism and retained for future reuse. Similarly, when prefilling and decoding are operated on the same processing unit, the newly generated response KV blocks may remain directly resident in HBM. Nevertheless, in either case, the alternative architecture can be accommodated by redefining the turn-wise increment $Z_j$ and the turn boundary.
\end{remark}


Our MCM is constructed under the following assumptions.

\begin{assumption}[Independent Block Increments]
    \label{ass:independent block incrememts}
    The block-increment pairs $\{(X_j,Y_j)\}_{j\geq1}$ (allowing arbitrary dependence across turns) are independent of the total number of turns $M$ within a conversation. The turn-wise content increments $Z_i$ satisfy that $Z_i \ge 1$ for all $i\ge 1$.
\end{assumption}

\begin{assumption}[Light-tailed Cumulative Block Workload]
    \label{ass: Regularity of one cumulative conversation blocks}
    The cumulative number of KV blocks generated over the lifetime of a conversation, defined as
    \begin{equation}
        \label{equ: cumulative number of blocks}
        B:=R_M=\sum_{\ell=1}^{M}Z_\ell,
    \end{equation}
    has a finite moment generating function in a positive neighborhood of $0$, i.e., there exists $\theta_0>0$ such that $\E[e^{\theta_0 B}]<\infty$.
\end{assumption}

\begin{assumption}[i.i.d. Prompt Interarrival Times with a Lipschitz-continuous CDF.]
    \label{ass: Exponential Prompt Inter-arrival Time}
    For each conversation, let $W_j$ denote the interarrival time between the submissions of the turn-$j$ prompt and the turn-$(j+1)$ prompt. The sequence $\{W_j\}_{j\geq1}$ is assumed to be i.i.d. with cumulative distribution function (CDF) $F$, independent of $M$ and the block increments $\{(X_j,Y_j)\}_{j\geq1}$. We further assume that $F$ is Lipschitz-continuous with constant $L_F$ and $\mu_F:= \E [W] <\infty$.
\end{assumption}



We note that the MCM reduces to the stochastic conversation model (SCM) introduced in \cite{zhang2026tail}, under the following additional assumptions:
\begin{enumerate}
    \item The block-increment pairs $\{(X_j,Y_j)\}_{j\geq1}$ are i.i.d.;
    \item The prompt interarrival times are i.i.d. exponentially distributed with rate $r>0$;
    \item The total number of turns $M$ is geometrically distributed. Specifically, after each completed turn, a conversation continues with probability $p\in(0,1)$ and terminates with probability $q:=1-p$. Equivalently, an active conversation can be viewed as having an independent new-prompt Poisson clock with rate $pr$ and an independent termination Poisson clock with rate $qr$.
\end{enumerate}
When the block increments additionally have constant means $\E[X_j]=k$ and $\E[Y_j]=m$ for all $j\geq1$, the expected history content size after turn $j$ can be simplified as
\begin{equation*}
    \bar R_j:=\E[R_j]=jk+(j-1)m.
\end{equation*}



\subsection{The  Hit Ratio}
\label{ssec: The  Hit Ratio}


We next introduce the hit ratio, which measures the cache reuse effectiveness of a management policy and will serve as the main performance metric throughout the paper. Unlike a request-level hit ratio that assigns equal weight to each cache access, the hit ratio quantifies the fraction of the overall KV block workload that can be directly served through history reuse from HBM. This block-level metric is particularly important for multi-turn LLM serving, because different conversations, and even different turns within the same conversation, may have substantially different history lengths.

Consider a generic conversation with $M$ turns and history content sizes $\{R_j\}_{j=1}^M$. If the conversation continues from turn $j$ to turn $j+1$, then the $R_j$ history blocks become reuse-eligible at the next prompt. Ideally, these blocks can be retrieved directly from HBM rather than reconstructed through reprefilling or transferred from lower-level storage such as CPU DRAM or SSD. Before decoding the turn-$(j+1)$ response, the previously generated history and the newly arrived turn-$(j+1)$ prompt must be organized as the complete prefix required by the model. We note that, on the one hand, recovering missing history from lower-level storage introduces additional communication and storage overhead. On the other hand, reprefilling the corresponding tokens incurs additional computation and increases the workload on the prefiller. In contrast, directly reusing resident KV blocks from HBM can reduce both data transfer overhead and recomputation cost. This motivates a workload-level metric that captures the amount of reusable KV history served from HBM.

We next specify the conditions under which a KV block stored in HBM can be successfully reused when its corresponding conversation receives a new prompt. In practical prefix-caching systems, two conditions are generally required:
\begin{enumerate}
    \item The block must be fully populated and therefore hashable, which means that its complete token content can be mapped to a deterministic identifier (a hash value) and stored in a hash map for later exact prefix matching and lookup. For instance, in our experiments (see \cref{sec: LLM-Inference Experiment and Empirical Results}), each block corresponds to $128$ tokens. A block having less than $128$ tokens becomes an unhashed block and cannot be reused. Such a phenomenon happens because the block is from the tail of a conversation and cannot fill a block.
    \item All preceding blocks in the same prefix must also be available\footnote{However, the preceding blocks do not necessarily need to reside in HBM. They may instead be retrieved from lower-level storage, such as CPU DRAM or SSD. After finding these prefix blocks, the subsequent blocks that remain resident in HBM can still be reused.}, because cache reuse follows the prefix-hit rule. For example, suppose a conversation contains blocks indexed from $0$ to $6$. Even if block-$5$ is fully populated and resident in HBM, it cannot be reused if block-$4$ is unavailable in any storage. In this case, block-$5$ must be reprefilled together with the missing preceding blocks.
\end{enumerate}


In our theoretical MCM analysis, we assume that all KV blocks are fully populated and hence hashable, so that condition 1 is always satisfied. The effect of unhashable tail blocks will be incorporated later in our practical MCM. We now define the hit ratio.


\begin{definition}
    \label{def:  hit ratio}
    For each turn $j\geq 2$, let the random variable $R^H_j \in \N$  denote the stationary number of history-prefix blocks reused directly from HBM under the LRU policy by a turn-$j$ access, with $0\leq R^H_j \leq R_{j-1}$. The \emph{hit ratio} under LRU is defined as
    \begin{equation}
    \label{equ:  hr}
        h :=\frac{\E \Big[\sum\limits_{j=2}^{M}R^H_{j} \Big]}{\E \Big[\sum\limits_{j=1}^{M} R_j\Big]}.
    \end{equation}
\end{definition}

In \eqref{equ:  hr}, the denominator $\E \Big[\sum\limits_{j=1}^{M} R_j\Big]$ represents the expected cumulative history content workload over the lifetime of a conversation, whereas the numerator represents the portion of this workload that is successfully touched through direct HBM reuse.


We use a simple example to illustrate this definition. Suppose a conversation contains three turns, with $X_1=2$, $X_2=4$, $Y_1=2$, $Y_2=5$, and $X_3=3$. Then the history content sizes are $R_1=X_1=2$, $R_2=X_1+(X_2+Y_1)=8$, and $R_3=X_1+(X_2+Y_1)+(X_3+Y_2)=16$. Suppose that, under LRU, all history blocks required at turns 2 and 3 are successfully reused from HBM. Then the  hit ratio is
\begin{equation*}
    h=\frac{0+2+8}{2+8+16}=\frac{10}{26}.
\end{equation*}

A key feature of the hit ratio is that its intrinsic upper bound is usually strictly smaller than $1$. This is because newly arrived prompt blocks cannot be reused from HBM, and only history contents associated with conversations that continue to the next turn are eligible for reuse. We therefore characterize the maximum upper bound on the hit ratio.

Define the sequence of tail probabilities 
\begin{equation}
    \label{equ: survival sequence}
    S_j:=\pro(M\geq j), \qquad j\geq 1,
\end{equation}
where $S_1=1$ and $S_j\downarrow 0$. For every $j$ such that $S_j>0$, define
\begin{equation}
    \label{equ: continuation probability at turn j}
    p_j:=\pro(M\geq j+1\mid M\geq j)=\frac{S_{j+1}}{S_j},
\end{equation}
as the continuation probability at turn-$j$, and $q_j:=1-p_j$ as the corresponding termination probability. The expected cumulative history content workload over a conversation lifetime is
\begin{equation}
    \label{equ: total number of blocks}
    A_R:= \mathbb{E} \left[ \sum\limits_{j=1}^M R_j \right ] =  \sum_{j\geq1}S_j\bar R_j,
\end{equation}
whereas the expected amount of reusable history content workload is
\begin{equation}
    \label{equ: number of blocks eligible for reuse}
    A_H:= \mathbb{E} \left[ \sum\limits_{j=1}^{M-1} R_j \right ] = \sum_{j\geq1}S_{j+1}\bar R_j.
\end{equation}
We define the intrinsic reuse potential of the workload as
\begin{equation}
    \label{equ: upper bound of all hit ratio}
    h_{\max}:=\frac{A_H}{A_R}.
\end{equation}
Then, for any HBM capacity,
\begin{equation*}
    h \leq h_{\max}.
\end{equation*}
The quantity $h_{\max}$ depends only on the statistical structure of the multi-turn conversation workload and is independent of the cache-management policy or the conversation-prompt arrival dynamics.


For instance, in our SCM model, where $p_j \equiv p$, $q_j \equiv 1-p$, and $\bar R_j = jk+(j-1)m$, we have that $A_R = \frac{k+pm}{q^2}$ and $A_H = \frac{p(k+pm)}{q^2}$. Hence, the intrinsic reuse potential has a closed form $h_{\max} = p$, which only depends on the continuation probability, or equivalently the mean number of turns in a conversation.



\subsection{The Least-Recently-Used (LRU) Policy}
\label{ssec: The LRU Policy}

In our MCM, we consider a whole-content version of the least-recently-used (LRU) policy. The system maintains an ordered resident list $\mathcal L$, from the most-recently-used (MRU) end to the LRU end. The list is updated whenever a new prompt is processed and the corresponding conversation content is inserted or refreshed in HBM. Suppose an arriving prompt belongs to a resident conversation $\theta\in\mathcal L$. The conversation label $\theta$ is removed from its current position and moved to the MRU end, while its previously stored content is replaced by the enlarged content containing the newly generated blocks. If the prompt belongs to a nonresident conversation $\theta\notin\mathcal L$, the conversation is inserted directly at the MRU end, together with its available history content and the newly generated blocks.


As specified in \cref{ssec: Generalized SCM}, the HBM can store at most $N$ KV blocks. Whenever inserting or updating a conversation content causes the total resident KV block workload to exceed this capacity, the system will evict complete conversation contents sequentially from the LRU end until the capacity constraint is restored.\footnote{In practical implementations, LRU eviction may operate at the block level and therefore allow partial-prefix eviction. We discuss these implementation differences in \cref{sec: LLM-Inference Experiment and Empirical Results}.} The whole-content LRU policy under our consideration is described by Algorithm \ref{alg: Whole-Content LRU} below. 

\begin{algorithm}[ht]
\caption{Whole-Content LRU}
\label{alg: Whole-Content LRU}
    \KwIn{HBM capacity $N$; ordered resident list $\mathcal L=(i_1,\ldots,i_K)$ from MRU to LRU, where $K$ is the number of conversations whose contents are stored in the HBM; resident content sizes $\{Z_i\}_{i\in\mathcal L}$; arriving turn $(\theta,j)$ with content size $R_j$.}
    
    \KwOut{Updated LRU list $\mathcal L$, content sizes $\{Z_i\}$, and hit indicator $H_{\theta,j}$.}
    
    $H_{\theta,j}\gets \one \{\theta\in\mathcal L\}$\;
    
    \If{$\theta\in\mathcal L$}{
        remove $\theta$ from its current position in $\mathcal L$\;
    }
    
    $Z_\theta\gets R_j$\;
    
    $\mathcal L\gets(\theta,\mathcal L)$
    \tcp*[r]{Move the complete updated content to MRU}
    
    \While{$\displaystyle\sum_{i\in\mathcal L} Z_i>N$}{
    
        $i^*\gets$ last element of $\mathcal L$
        \tcp*[r]{Current LRU conversation}
    
        remove $i^*$ from $\mathcal L$\;
    
        delete $Z_{i^*}$
        \tcp*[r]{Whole-content eviction}
    }
    \Return{$(\mathcal L,\{Z_i\},H_{\theta,j})$}\;
\end{algorithm}

Under the whole-content LRU policy in our MCM, a conversation content is either fully resident in HBM or entirely absent. If the content is missed when its next-turn prompt arrives, its history blocks must be retrieved from lower-level storage or reconstructed through reprefilling. Therefore, the  hit ratio admits the following equivalent characterization.

For each turn $j\geq 2$, define the LRU hit indicator
\[H_{j} := \one \left\{ \text{the history content of size $R_{j-1}$ is resident in HBM when turn $j$ arrives} \right\}.
\]
If $H_j=1$, all $R_{j-1}$ history blocks can be directly reused from HBM. Otherwise, none of these blocks is touched through HBM reuse under the whole-content convention. Hence, the  hit ratio in \eqref{equ:  hr} can be equivalently written as
\begin{equation}
    \label{equ:  hr LRU}
    h  :=\frac{\E \Big[\sum\limits_{j=2}^{M}R_{j-1}H_{j}\Big]}{\E \Big[\sum\limits_{j=1}^{M} R_j\Big]}.
\end{equation}


\subsection{The Mean-Field Asymptotic Regime}
\label{ssec: Our Mean-Field Model}


We next introduce the mean-field scaling used in our theoretical analysis. The motivation is that, if the HBM capacity $N$ is very large while the conversation arrival process remains unchanged, an arriving prompt can retain almost all reusable history blocks in HBM, and the hit ratio approaches the upper bound $h_{\max}$. Similarly, when $N$ is very small, only a negligible fraction of the history workload can remain resident in HBM, and the hit ratio approaches $0$. To obtain a nontrivial asymptotic regime, we therefore scale the conversation arrival rate proportionally with the HBM capacity.

Specifically, for the system indexed by $N$, we set
\begin{equation*}
    \lambda_0^{(N)}:=N\lambda_0,
\end{equation*}
while keeping the prompt interarrival distribution $F$, the turn-number distribution $M$, and the block-increment sequence $\{(X_j,Y_j)\}_{j\geq1}$ unchanged. Under this scaling, both the HBM capacity and the aggregate workload arrival rate grow linearly with $N$, and we will show in \cref{ssec: Mean-Field Results of the MCM under LRU} that this scaling yields a nondegenerate mean-field limit for the cache dynamics and hit ratio.


\begin{remark}
    \label{rmk: different scaling way}

    The mean-field scaling above has an equivalent interpretation obtained through a rescaling of prompt interarrival time. Specifically, suppose that we keep the conversation arrival rate fixed at $\lambda_0$, while scaling the prompt interarrival distribution as $F^{(N)}(\cdot):=F(\cdot/N)$, i.e., the original interarrival times are scaled by a factor of $N$. Under the scaled time $t:=t/N$, conversations then arrive at rate $N\lambda_0$, while the interarrival distribution becomes $F$, which is exactly the mean-field scaling considered above.

    This alternative interpretation is closer to our LLM inference experiments. In practice, increasing the physical conversation arrival rate proportionally with $N$ may overload the prefill and decode workers and make computation be the system bottleneck. The resulting end-to-end (E2E) latency may then substantially affect the observed prompt interarrival times. Therefore, we use the arrival-rate scaling $\lambda_0^{(N)}=N\lambda_0$ for the theoretical mean-field analysis. Nonetheless, in the experiments we keep the physical conversation arrival rate at a practical level and adjust the prompt interarrival times to emulate the corresponding mean-field regime.
\end{remark}

\subsection{Tagged-conversation Displacement Process and Characteristic Time}
\label{ssec: tagged-conversation Displacement Process and Characteristic Time}

Our mean-field analysis uses a classical tagged-job approach in queueing theory. In our setting, we tag a conversation immediately after its turn-$j$ access and study its subsequent evolution in stationarity. We introduce two key notions: the tagged-conversation displacement process and the corresponding characteristic time. These objects provide the basis for our mean-field analysis of eviction and the LRU hit ratio.

We first formulate a tagged-conversation representation of the whole-content LRU dynamics. The key idea is to characterize the eviction of a tagged content through the cumulative KV block mass that moves ahead of it, after its most recent access. This transforms the evolution of the LRU ordering into a one-dimensional cumulative displacement process.

We tag a conversation immediately after one of its prompt accesses. Suppose that the tagged conversation receives its turn-$j$ prompt at time $0$ and is moved to the MRU end of the whole-content LRU list. At this moment, no other resident blocks lie ahead of the tagged content in the LRU ordering. Assuming that the tagged conversation receives no further access before eviction, as other conversations are subsequently accessed, their KV blocks move to the MRU side and therefore accumulate in front of the tagged content.

A background conversation that has already moved ahead of the tagged content should not contribute its entire content again at its later accesses. Once its pre-existing history has crossed the tagged content, only the newly appended KV blocks generated at subsequent accesses further increase the block mass ahead of the tagged content. Hence, for each background conversation $i$, define its cumulative displacement contribution $C_i(t)$ as follows:
\begin{enumerate}
    \item Before its first access after time $0$, $C_i(t)=0$;
    \item At its first access after time $0$, its complete pre-existing history content and the newly accessed prompt's blocks cross the tagged content;
    \item At each subsequent access, only the newly appended KV block mass contributes additional displacement.
\end{enumerate}
We are now ready to define the tagged-conversation displacement process.


\begin{definition}[(Tagged Turn-$j$ Conversation) Displacement Process]
    \label{def: tagged-conversation displacement process}
    In the MCM-LRU system with HBM capacity $N$, consider a tagged conversation immediately after its turn-$j$ access at time $0$, with history content size $R_j$.
    We condition on the event that its content is evicted before its next reuse. The displacement process of this content immediately after its turn-$j$ access is defined as
    \begin{equation}
    \label{equ: tagged displacement process}
        D_{N,j}(t):=\sum_{i\neq \mathrm{tag}} C_i(t), \qquad t\geq0.
    \end{equation}
    The corresponding normalized displacement process is defined as $\frac{D_{N,j}(t)}{N}$.
\end{definition}

Under the whole-content LRU convention, $D_{N,j}(t)$ is nondecreasing in $t$. We next define the corresponding characteristic time, denoted by $\tau_C^{(N)}(j)$.


\begin{definition}[Tagged Turn-$j$ Characteristic Time]
    \label{def: characteristic time}
    Consider a tagged conversation immediately after its turn-$j$ access at time $0$, with history content size $R_j$. We condition on the event that its content is evicted before its next reuse. The characteristic time $\tau_C^{(N)}(j)$ is defined as the time required for the background displacement to evict this tagged content from HBM, assuming that the tagged conversation is not accessed again before eviction. Specifically,
    \begin{equation}
    \label{equ: exact characteristic time}
        \tau_C^{(N)}(j):=
        \begin{cases}
            0, & R_j>N, \\
            \inf\{t\geq 0: D_{N,j}(t)>N-R_j\}, & R_j\leq N.
        \end{cases}
    \end{equation}
\end{definition}

In \cref{ssec: Mean-Field Results of the MCM under LRU}, we establish the convergence of the normalized displacement process $\frac{D_{N,j}(t)}{N}$ and the characteristic time $\tau_C^{(N)}(j)$, and then use these results to derive the mean-field convergence of the LRU hit ratio.


\section{Extended mean-field results}
\label{sec: main results}

In this section, we present the main theoretical results of our mean-field analysis. We first establish the convergence of the displacement process, the characteristic time, and the LRU hit ratio for the multi-turn conversation model (MCM) in \cref{ssec: Mean-Field Results of the MCM under LRU}. In \cref{ssec: Practical MCM Consequence and Theoretically-Motivated Hit Ratio Estimator}, we extend the analysis to a practical MCM that incorporates additional features of real multi-turn LLM serving systems. Based on its mean-field convergence results, we develop a theoretically-motivated estimator for the  hit ratio. The accuracy of this estimator will be evaluated empirically in \cref{sec: LLM-Inference Experiment and Empirical Results}. Unless otherwise stated, all results in this section concern a stationary system.


\subsection{Mean-Field Results of the MCM under LRU}
\label{ssec: Mean-Field Results of the MCM under LRU}


We begin with the mean-field limit of the (tagged-conversation) displacement process. Its normalized sample path converges uniformly on every finite time horizon to the deterministic mean displacement curve. We first define this mean displacement curve.

\begin{definition}
    \label{def: Mean displacement Curve}
    Recall that $A_H:=\sum_{j\geq1}S_{j+1}\bar R_j$ denotes the expected amount of history content workload eligible for reuse, and that $B:=R_M$ denotes the total KV block count required by the end of a conversation. The mean displacement curve $d:[0,\infty)\rightarrow\mathbb R$ is defined as
    \begin{equation}
    \label{equ: Mean Displacement Curve-1}
        d(t):=\lambda_0\E[B]t+\lambda_0 A_H\int_0^t\bar F(u)\,du.
    \end{equation}
\end{definition}



\begin{theorem}[Sample-Path Convergence of the Normalized Displacement Process]
\label{thm: Displacement Process Convergence}
    The tagged turn-$j$ content process follows the same law, such that $D_{N,j}(\cdot)\overset{d}{=}D_N(\cdot)$ for every turn state $j$. In addition, for a fixed $T<\infty$, let $V_T  = \lambda_0T\E[B^2]+\lambda_0\mu_W\E[(M-1)B^2]<\infty.$, and let $L :=\sup_{0\leq t\leq T}d'(t)=\lambda_0\E[B]+\lambda_0A_H$.
    Then, for every $t\in[0,T]$ and every $\varepsilon>0$,
    \begin{equation}
    \label{equ: displacement pointwise concentration}
        \pro\left(\left|\frac{D_N(t)}{N}-d(t)\right|>\varepsilon\right)\leq\frac{V_T}{N\varepsilon^2}.
    \end{equation}
    Moreover, for every $0<\varepsilon\leq1$,
    \begin{equation}
    \label{equ: displacement sample path concentration}
        \pro\left(\sup_{0\leq t\leq T}\left|\frac{D_N(t)}{N}-d(t)\right|>\varepsilon\right)\leq\frac{C_T}{N\varepsilon^3},
    \end{equation}
    where $C_T:=\frac{16V_T}{9}(2+4 L T).$
    
    Consequently, for any fixed $T\geq 0$,
    \begin{equation}
    \label{equ: displacement sample path convergence}
        \sup_{0\leq t\leq T}\left|\frac{D_N(t)}{N}-d(t)\right|\xrightarrow{\mathbb P}0.
    \end{equation}
    In particular, the finite-size bounds imply
    \begin{equation}
    \label{equ: displacement pointwise rate}
        \left|\frac{D_N(t)}{N}-d(t)\right|=O_{\mathbb P}(N^{-1/2})
    \end{equation}
    for every fixed $t$, and
    \begin{equation}
    \label{equ: displacement uniform rate}
        \sup_{0\leq t\leq T}\left|\frac{D_N(t)}{N}-d(t)\right|=O_{\mathbb P}(N^{-1/3}).
    \end{equation}
\end{theorem}

We defer the proof of \cref{thm: Displacement Process Convergence} to \cref{ssec: Proof of the full version of the sample path convergence}. Nonetheless, we provide a proof sketch.

\noindent\textit{Proof Sketch:} We first represent the stationary background conversations as a marked Poisson random measure. Hence, we can express the displacement process as an integral of their cumulative displacement contributions. This representation also shows that the law of the background displacement process is invariant w.r.t. the tagged turn state $j$. We then derive the mean displacement curve $d(t)$ by separately evaluating the expected contributions from pre-existing history contents and newly generated KV blocks. The Poisson random measure representation further yields a pointwise second-moment bound of order $O(N^{-1})$ for the normalized displacement process. Finally, using the nondecreasing property of $D_N(t)$ and $d(t)$, the Lipschitz continuity of $d(t)$, and a discretization argument over a finite time horizon with granularity $O(N^{\frac{1}{3}})$, we obtain the uniform convergence. \hfill $\square$


The sample-path convergence of the displacement process leads to our next core result, the tagged turn-$j$ content characteristic time converges to a deterministic limit. We first define this limiting mean-field characteristic time. Recall from \cref{def: Mean displacement Curve} that $d(0)=0$ and $d(t)\rightarrow\infty$ as $t\rightarrow\infty$. Moreover, $d$ is differentiable with uniformly lower bounded derivative 
\begin{equation*}
    d'(t)=\lambda_0\E[B]+\lambda_0A_H\bar F(t) \geq \lambda_0\E[B]>0.
\end{equation*}
Hence, $d$ is strictly increasing, and its inverse $d^{-1}:[0,\infty)\rightarrow[0,\infty)$ is well-defined.
\begin{definition}
    \label{def: deterministic mean-field characteristic time}
    The \emph{mean-field characteristic time} $T_C$ is defined as $T_C:=d^{-1}(1)$, or equivalently, as the unique solution to the equation
    \begin{equation}
        \label{equ: mean displacement curve solution to 1}
        d(T_C)=\lambda_0\E[B]T_C+\lambda_0A_H\int_0^{T_C}\bar F(u)\,du=1.
    \end{equation}
\end{definition}


We now state our second core result on the convergence of the characteristic time.

\begin{theorem}[Characteristic Time Convergence]
\label{thm: Characteristic Time Convergence- Full}
    Fix a turn state $j$ such that $S_j>0$, and a finite horizon $T>T_C$. For every
    \begin{equation*}
        0<\varepsilon<\min\{T_C,T-T_C\},
    \end{equation*}
    the characteristic time satisfies
    \begin{equation}
        \label{equ: characteristic time finite N bound}
        \pro\left(\left|\tau_C^{(N)}(j)-T_C\right|>\varepsilon\right)
        \leq\frac{18V_T}{Nc_d^2\varepsilon^2}
        +\frac{\E[e^{\theta B}]}{S_j}\exp\left(-\frac{\theta c_d}{3}N\varepsilon\right),
    \end{equation}
    where $c_d:=\lambda_0\E[B]$ and $V_T:=\int_{\mathcal Z}G_T(z)^2\,\nu(dz)<\infty$ is the second-moment constant introduced in \cref{lem: Poisson Random Measure Representation}. Consequently,
    \begin{equation}
    \label{equ: characteristic time convergence-2}
        \tau_C^{(N)}(j)\xrightarrow{\mathbb P}T_C
    \end{equation}
    for every fixed turn state $j$ with $S_j\geq 0$, and
    \begin{equation*}
        \tau_C^{(N)}(j)-T_C=O_{\mathbb P}(N^{-1/2}).
    \end{equation*}
\end{theorem}

The proof of \cref{thm: Characteristic Time Convergence- Full} is deferred to \cref{ssec: Proof of the Full Version of CT convergence}. We also provide a proof sketch here.

\noindent\textit{Proof Sketch:} We note that different tagged turn-$j$ contents have different eviction thresholds, because eviction occurs when $D_{N,j}(t)>N-R_j$. However, for any fixed $j$ with $S_j>0$, the content size satisfies $R_j=O_{\mathbb P}(1)$, and hence $R_j/N$ is asymptotically negligible. To control the event $\{|\tau_C^{(N)}(j)-T_C|>\varepsilon\}$, it is sufficient to control the normalized displacement process at the two deterministic times $T_C-\varepsilon$ and $T_C+\varepsilon$, and the deviation of $R_j/N$ from $0$. Because $d(T_C-\varepsilon)<1<d(T_C+\varepsilon)$, the pointwise concentration bounds for $|\frac{D_N(t)}{N}-d(t)|$ imply that, with high probability, the tagged content is resident at time $T_C-\varepsilon$ and already evicted by time $T_C+\varepsilon$. By applying a union bound to these deviations, we obtain the convergence of $\tau_C^{(N)}(j)$ to $T_C$. The $O_{\mathbb P}(N^{-1/2})$ convergence rate follows from the pointwise concentration rate of the displacement process in \cref{thm: Displacement Process Convergence}. \hfill $\square$


We finally turn to the third core result for the MCM-LRU system: the convergence of the  hit ratio. In the system indexed by $N$, for each turn $j\geq2$, define the hit indicator
\[
    H_{N,j} := \one \left\{ \text{the history content of size $R_{j-1}$ is resident in HBM when turn $j$ arrives} \right\}.
\]
Then the finite-$N$  hit ratio under whole-content LRU can be written as
\begin{equation}
\label{equ: finite N LRU hit ratio}
    h ^{(N)}:=\frac{\E\left[\sum_{j=1}^{M-1}R_jH_{N,j+1}\right]}{A_R},
\end{equation}
where we recall that $A_R=\sum_{j\geq1}S_j\bar R_j.$ We define the corresponding mean-field limiting hit ratio as
\begin{equation}
\label{equ: limiting LRU hit ratio}
    h ^{(\infty)}  :=\frac{A_H}{A_R}  F(T_C). 
\end{equation}
We now state \cref{thm: Hit Ratio Convergence Full}, which follows from the convergence of the displacement process and the characteristic time established above. Its proof is deferred to \cref{ssec: Proof of the Full Version of hit ratio conv MCM-LRU short}.


\begin{theorem}[Hit Ratio Convergence in the MCM-LRU Model]
\label{thm: Hit Ratio Convergence Full}
    Under the MCM assumptions, fix a finite $J$ such that $S_{J+1}>0$, and fix $T>T_C$. For every $0<\varepsilon<\min\{T_C,T-T_C\}$,
    \begin{equation}
    \label{equ: finite N hit ratio error bound}
    \left|h_{\LRU}^{(N)} -h_{\LRU}^{(\infty)}  \right| \leq \frac{L_F\varepsilon A_H^{(J)}}{A_R} +\frac{K_J}{A_R}\sqrt{p_{N,J}(\varepsilon)} +\frac{2\E[B^2e^{\theta B}]e^{-\theta(J+2)}}{A_R},
    \end{equation}
    where $p_{N,j}(\varepsilon):= \frac{18V_T}{Nc_d^2\varepsilon^2}
        +\frac{\E[e^{\theta B}]}{S_j}\exp\left(-\frac{\theta c_d}{3}N\varepsilon\right)$ and  $K_J:=\sum_{j=1}^{J}S_{j+1}\sqrt{\E[R_j^2]}$.
    Consequently, the system-level hit ratio under the whole-content LRU policy converges when $N\rightarrow\infty$:
    \begin{equation}
    \label{equ: hit ratio convergence-2}
        h_{\LRU}^{(N)} \rightarrow h_{\LRU}^{(\infty)}  =\frac{A_H}{A_R} F(T_C).
    \end{equation}
\end{theorem}


\noindent\textit{Proof Sketch:} The proof mainly relies on the concentration of the characteristic time established in \cref{thm: Characteristic Time Convergence- Full} and the Lipschitz continuity of the interarrival CDF $F$. For each fixed turn state $j$, \cref{thm: Characteristic Time Convergence- Full} suggests that its hit probability is asymptotically characterized by $F(T_C)$. However, the characteristic time convergence is not uniform over all turn states. We therefore truncate the reusable-workload series at some level $J$, and compare $\sum_{j=1}^{J}S_{j+1}\bar R_jF(T_C)$ with $\E\big[\sum\limits_{j=1}^{\min\{M-1,J\}}R_jH_{N,j+1}\big].$ For every fixed $J$, the difference between these two quantities vanishes as $N\rightarrow\infty$ by the characteristic time convergence and the Lipschitz continuity of $F$. The remaining contribution from turn states $j>J$ is controlled by the tail reusable workload $\sum_{j>J}S_{j+1}\bar R_j$, which vanishes as $J\rightarrow\infty$. Finally, we conclude \cref{thm: Hit Ratio Convergence Full} by first taking $N\rightarrow\infty$ and then $J\rightarrow\infty$. \hfill $\square$


\subsection{Practical MCM Consequence and Theoretically-Motivated Hit Ratio Estimator}
\label{ssec: Practical MCM Consequence and Theoretically-Motivated Hit Ratio Estimator}


In real multi-turn LLM conversation systems, the last KV block generated at each turn is usually partially filled. Therefore, this block does not have a valid hash. We call this block the \emph{unhashable tail block}. To capture this feature, we introduce a practical multi-turn conversation model (practical MCM). As a modeling approximation, we conservatively assume each completed turn generates one unhashable tail block. Although this block occupies one physical HBM block, it cannot participate in prefix reuse at subsequent turns and can only gradually move toward the LRU end until eviction.

When the conversation proceeds to the next turn, the tokens contained in the unhashable tail block of turn $j$ must be reprefilled together with the turn-$j$ response tokens and the turn-$(j+1)$ prompt tokens. Accordingly, these tokens can be absorbed into the next turn-wise content increment. % Without loss of generality, we therefore write $Z_1:=X_1$ and $Z_{j+1}:=Y_j+X_{j+1}$ for $j\geq1$, where $Z_{j+1}$ includes the tokens carried by the unhashable tail block from turn $j$.



\subsubsection{Practical Characteristic Time and Hit Ratio with Unhashable Tail Blocks}
\label{sssec: Practical Hit Ratio}


We keep the definitions of $R_j$ and $B$ as the KV block counts, including the unhashable tail blocks. Hence, the total number of newly generated physical KV blocks is unchanged from the original MCM. The only modification is that, when a conversation is revisited after turn $j$, the unhashable tail block cannot participate in prefix reuse, so only $R_j-1$ history blocks are hashable and are moved back to the MRU side. We therefore define the expected reusable hashable workload as
\begin{equation}
\label{equ: practical hashed AH}
    A_H^{\mathrm{hash}}:=\sum_{j\geq1}S_{j+1}\E[R_j-1]=A_H-\E[M-1].
\end{equation}


Let $D_N^{\mathrm{prac}}(t)$ denote the displacement process in the practical system. Because every newly generated physical KV block still contributes to displacement, the new-block contribution remains unchanged. In contrast, each revisit moves one fewer pre-existing history block because the unhashable tail block cannot participate in prefix reuse. Thus, we define the practical mean displacement curve as
\begin{equation}
\label{equ: practical displacement curve}
    d_{\mathrm{prac}}(t):=\lambda_0 \E[B]t+\lambda_0 A_H^{\mathrm{hash}}\int_0^t \bar F(u) du.
\end{equation}
Because $d_{\mathrm{prac}}(0)=0$, $\lim\limits_{t\uparrow \infty} d_{\mathrm{prac}}(t)=\infty$, and
\begin{equation*}
    d_{\mathrm{prac}}'(t)=\lambda_0\E[B]+\lambda_0 A_H^{\mathrm{hash}} \bar F(t)>0,
\end{equation*}
there exists a unique practical mean-field characteristic time $T_C^{\mathrm{prac}}$ satisfying
\begin{equation}
\label{equ: practical characteristic time}
    d_{\mathrm{prac}}(T_C^{\mathrm{prac}})=1.
\end{equation}

For a turn-$j$ tagged conversation, let $\tau_{C,\mathrm{prac}}^{(N)}(j)$ denote the HBM residence time of its hashable history content with size $R_j-1$, in the absence of a next-turn access. %We further let $\tau_{\mathrm{nohash}}^{(N)}$ denote the residence time of a tagged unhashable block, and let $N_{\mathrm{nohash}}^{(N)}$ denote the number of unhashable blocks resident in HBM at stationarity.

Finally, we define the practical finite-$N$ hit ratio as
\begin{equation}
\label{equ: practical finite N hit ratio}
    h_{ \mathrm{prac}}^{(N)}:=\frac{\E\left[\sum_{j=1}^{M-1}(R_j-1)H_{N,j+1}^{\mathrm{prac}}\right]}{A_R},
\end{equation}
where $H_{N,j+1}^{\mathrm{prac}}$ is the indicator of the event that the hashable history content after turn $j$ remains resident in HBM when the turn-$(j+1)$ prompt arrives. The corresponding mean-field limiting hit ratio is defined as
\begin{equation}
\label{equ: practical limiting hit ratio}
    h_{ \mathrm{prac}}^{(\infty)}:=\frac{A_H^{\mathrm{hash}}}{A_R}F(T_C^{\mathrm{prac}}).
\end{equation}


We next state the main mean-field convergence result for the practical MCM-LRU system.

\begin{theorem}[Practical Hit Ratio Convergence]
\label{thm: Practical Hit Ratio Convergence Full}
    Fix a finite $J$ such that $S_{J+1}>0$ and $T>T_C^{\mathrm{prac}}$. For every $0<\varepsilon<\min\{T_C^{\mathrm{prac}},T-T_C^{\mathrm{prac}}\},$ the difference between the practical finite-size hit ratio and the mean-field limiting hit ratio can be bounded, such that
    \begin{equation}
    \label{equ: practical finite N hit ratio bound}
        \left|h_{ \mathrm{prac}}^{(N)}-h_{ \mathrm{prac}}^{(\infty)}\right| \leq \frac{L_F\varepsilon A_{H,\mathrm{hash}}^{(J)}}{A_R} + \frac{K_{J,\mathrm{hash}}}{A_R}\sqrt{p_{N,J}^{\mathrm{prac}}(\varepsilon)} + \frac{2\E[B^2e^{\theta B}]e^{-\theta(J+2)}}{A_R},
    \end{equation}
    where $L_F$ is the Lipschitz constant of $F$, $p_{N,J}^{\mathrm{prac}}(\varepsilon) :=\frac{18V_T^{\mathrm{prac}}}{Nm_{\mathrm{prac}}^2\varepsilon^2} + \frac{\E[e^{\theta B}]}{S_J} \exp\left(-\frac{\theta m_{\mathrm{prac}}}{3}N\varepsilon\right)$, and $ K_{J,\mathrm{hash}}:=\sum_{j=1}^{J}S_{j+1}\sqrt{\E[(R_j-1)^2]}$.  Consequently,
    \begin{equation}
    \label{equ: practical hit ratio convergence full}
        h_{ \mathrm{prac}}^{(N)}\rightarrow h_{ \mathrm{prac}}^{(\infty)}= \frac{A_H^{\mathrm{hash}}}{A_R} F(T_C^{\mathrm{prac}}).
    \end{equation}
\end{theorem}

The proof of \cref{thm: Practical Hit Ratio Convergence Full} is deferred to \cref{ssec: Proof of the Full Version of practical results}. It follows the same sequence as in the original MCM-LRU analysis, from displacement-process convergence, via characteristic time convergence, to hit-ratio convergence, which is the same as the map of proofs of \cref{thm: Displacement Process Convergence,thm: Characteristic Time Convergence- Full,thm: Hit Ratio Convergence Full}.


\subsubsection{Theoretically-Motivated Hit Ratio Estimator}
\label{sssec: Theoretically-Motivated Hit Ratio Estimator}

Inspired by the mean-field convergence results in \cref{thm: Practical Hit Ratio Convergence Full}, we develop a theoretically-motivated estimator for practical LLM inference systems.

Suppose that the empirical workload contains $L$ conversations. For conversation $i$, let $M_i$ denote its number of turns, $R_{i,j}$ the physical history content size after turn $j$, and $B_i:=R_{i,M_i}$ its total lifetime KV block count. We define the empirical mean cumulative history content workload as
\begin{equation}
\label{equ: practical estimator AR}
    \widehat A_R := \frac{1}{L} \sum_{i=1}^{L} \sum_{j=1}^{M_i}R_{i,j},
\end{equation}
and the empirical mean hashable workload as
\begin{equation}
\label{equ: practical estimator AH hash}
    \widehat A_H^{\mathrm{hash}} := \frac{1}{L} \sum_{i=1}^{L} \sum_{j=1}^{M_i-1}(R_{i,j}-1).
\end{equation}
The empirical mean total lifetime KV block count and the empirical mean number of turns are defined, respectively, as
\begin{equation}
\label{equ: practical estimator moments}
    \widehat{\E[B]} := \frac{1}{L}\sum_{i=1}^{L}B_i, \quad \text{and}\quad \widehat{\E[M]} :=\frac{1}{L}\sum_{i=1}^{L}M_i.
\end{equation}

Let $\widehat F$ denote the empirical CDF of the prompt interarrival time, and let $\widehat{\bar F}(t):=1-\widehat F(t)$ denote its tail. The practical mean-field characteristic time estimator $\widehat T_C^{\mathrm{prac}}$ is defined as the unique positive solution to the equation
\begin{equation}
\label{equ: practical TC estimator}
    \lambda_0 \widehat{\E[B]}\widehat T_C^{\mathrm{prac}}+\lambda_0 \widehat A_H^{\mathrm{hash}}\int_0^{\widehat T_C^{\mathrm{prac}}}\widehat{\bar F}(u)\,du=N.
\end{equation}
The corresponding mean-field estimator of the practical hit ratio is
\begin{equation}
\label{equ: practical hit ratio estimator}
    \widehat h_{ \mathrm{prac}}^{\mathrm{MF}}:=\frac{\widehat A_H^{\mathrm{hash}}}{\widehat A_R}\widehat F\left(\widehat T_C^{\mathrm{prac}}\right).
\end{equation}


In practice, the complete conversation-level workload trace may not be fully available or predictable. Nonetheless, our estimator does not require direct access to every individual conversation, as long as reasonable prior or externally estimated values of the required workload statistics and the prompt interarrival distribution are available. Specifically, we can substitute the estimates of $\E[B]$, $A_R$, $A_H^{\mathrm{hash}}$, and $F$ into the same mean-field formulas. Therefore, this theoretically-motivated hit ratio estimator remains applicable when partial workload distribution and system dynamics information is available.


\section{Proofs}
\label{sec: Proofs}

In this section, we provide the proofs of our main results in \cref{sec: main results}. Our main MCM-LRU mean-field result, the normalized tagged-$j$ conversation displacement process convergence (\cref{thm: Displacement Process Convergence}), the tagged-$j$ characteristic time convergence (\cref{thm: Characteristic Time Convergence- Full}), and the hit ratio convergence (\cref{thm: Hit Ratio Convergence Full}), are demonstrated in \cref{ssec: Proof of the full version of the sample path convergence,ssec: Proof of the Full Version of CT convergence,ssec: Proof of the Full Version of hit ratio conv MCM-LRU short}, respectively. 

The proof of \cref{thm: Practical Hit Ratio Convergence Full} follows essentially the same arguments as those developed for the MCM-LRU system. We will demonstrate it in \cref{ssec: Proof of the Full Version of practical results}.


\subsection{Proof of \cref{thm: Displacement Process Convergence}}
\label{ssec: Proof of the full version of the sample path convergence}


We begin by demonstrating the sample-path convergence of the tagged-conversation displacement process. The proof proceeds in three steps. In \cref{sssec: Representation of the Displacement Process}, we first establish a Poisson random measure representation of the displacement process and show in \cref{lem: Poisson Random Measure Representation}, that $D_{N,j}(\cdot)\stackrel{\mathcal L}{=}D_N(\cdot)$ for every turn state $j$. In \cref{sssec: Displacement Process Decomposition}, we decompose the displacement into two subprocesses: the pre-existing history content contribution and the newly generated KV block contribution. Their respective mean curves are derived in \cref{lem: Initial-prefix contribution,lem: Mean New Blocks Contribution} respectively, which sum to the mean displacement curve $d(t)$ in \cref{lem: Mean Displacement Curve}. Finally, in \cref{sssec: normalized displacement sample path convergence}, we prove \cref{thm: Displacement Process Convergence}, by combining the results in \cref{lem: Poisson Random Measure Representation,lem: Mean Displacement Curve}.


\subsubsection{Poisson Random Measure Representation of the Displacement Process}
\label{sssec: Representation of the Displacement Process}

We first establish a Poisson random measure representation of the tagged-conversation displacement process in \cref{lem: Poisson Random Measure Representation}. Using this Poisson random measure representation, we are prepared to then analyze its mean curve, and finally show the sample path convergence of our normalized displacement process.

\begin{lemma}[Poisson Random Measure Representation of the Displacement Process]
\label{lem: Poisson Random Measure Representation}
    Fix a finite horizon $T<\infty$. Recall \cref{ass: Exponential Prompt Inter-arrival Time}, where $\{W_j\}_{j\geq1}$ are i.i.d. interarrival times with cumulative distribution function $F$ and finite mean $\mu_W:=\E[W_1]$, independent of the conversation mark $(M,\{X_j,Y_j\}_{j\geq1})$. Then, for every fixed tagged turn state $j$ with $S_j>0$, the displacement process $\{D_{N,j}(t):0\leq t\leq T\}$ admits the representation
    \begin{equation}
    \label{equ: Poisson Random Measure Representation}
        D_{N,j}(t)\stackrel{\mathcal L}{=}D_N(t):=\int_{\mathcal Z}\Phi_t(z)\,\mathcal N_N(dz), \qquad 0\leq t\leq T,
    \end{equation}
    where $\mathcal N_N$ is a Poisson random measure on the marked conversation trajectory space $\mathcal Z$ with intensity measure $N\nu$. The path functional $\Phi_t(z)$ denotes the cumulative KV block mass contributed by the background conversation $z$ to the displacement process during $(0,t]$.

    For a marked conversation $z=(a,\zeta)$, where $a$ denote its arrival time and $\zeta$ denote its mark, we define its lifetime span $L(\zeta):=\sum_{\ell=1}^{M-1}W_\ell,$ and the envelope
    \begin{equation*}
        G_T(a,\zeta):=B(\zeta)\one\{-L(\zeta)<a\leq T\}.
    \end{equation*}
    Then, for every $z\in\mathcal Z$,
    \begin{equation}
    \label{equ: PRM contribution bound}
        0\leq\Phi_s(z)\leq\Phi_t(z)\leq G_T(z), \qquad 0\leq s\leq t\leq T.
    \end{equation}
    Moreover,
    \begin{equation}
    \label{equ: PRM second moment constant}
        V_T:=\int_{\mathcal Z}G_T(z)^2\,\nu(dz)=\lambda_0T\E[B^2]+\lambda_0\mu_W\E[(M-1)B^2]<\infty.
    \end{equation}
    % Consequently, the law of the background displacement process is invariant with respect to the tagged turn state $j$.
\end{lemma}



\begin{proof}
We first construct the conversation arrivals on the entire time axis $\mathbb R$, which is a Poisson process with rate $N\lambda_0$. Each arrival at time $a\in\mathbb R$ is independently assigned a mark
\begin{equation*}
    \zeta=\left(M,\{X_\ell,Y_\ell\}_{\ell\geq1},\{W_\ell\}_{\ell\geq1}\right)\in\mathcal M,
\end{equation*}
which specifies the complete trajectory of the corresponding conversation in the mark space $\mathcal M$. We then let $P_\zeta$ denote the distribution of the conversation mark $\zeta$, and let $z=(a,\zeta)$ additionally record the conversation arrival time. The collection of marked conversation trajectories then forms a Poisson random measure $\mathcal N_N$ on $\mathcal Z:=\mathbb R\times\mathcal M$, with intensity measure
\begin{equation*}
    N\nu(dz)=N\lambda_0\,da\,P_\zeta(d\zeta).
\end{equation*}

For a marked background conversation $z=(a,\zeta)$, we let $\Phi_t(z)$ denote the total KV block mass contributed by this conversation that moves ahead of the tagged content during $(0,t]$.  Under the Palm construction for the tagged conversation, Slivnyak's theorem (see, e.g., Theorem 1.4.5 in \cite{baccelli2010stochastic}) implies that, after removing the tagged conversation, the remaining marked background conversation trajectories have the same law as the original stationary marked Poisson process. Moreover, the trajectory of the tagged conversation is independent of the background process. Because the background marked Poisson process is stationary, the displacement process observed from the tagged turn-$j$ access at time $t=0$ has the same law for every fixed $j$ with $S_j>0$.

If a background conversation is already active at time $0$, then at its next prompt access, its pre-existing history content crosses in front of the tagged content. After this access, those existing blocks remain ahead of the tagged content, so subsequent accesses contribute only newly generated blocks. If a background conversation arrives after time $0$, all of its blocks generated by time $t$ contribute in the same manner. Hence every KV block of a background conversation contributes to displacement at most once. 

For a conversation mark $\zeta$, define its lifetime span by $L(\zeta):=\sum_{\ell=1}^{M-1}W_\ell.$ A background conversation can contribute to the displacement over the horizon $[0,T]$ only if its arrival time $a$ satisfies $a\leq T$ and $a+L(\zeta)>0$, or equivalently, $-L(\zeta)<a\leq T$. Therefore, for every $z=(a,\zeta)\in\mathcal Z$, the path $\Phi_\cdot(z)$ is nondecreasing and satisfies that
\begin{equation*}
    0\leq\Phi_s(z)\leq\Phi_t(z)\leq B(\zeta)\one\{-L(\zeta)<a\leq T\}=G_T(z), \qquad 0\leq s\leq t\leq T.
\end{equation*}
We then sum the contributions over all background conversations, and have that
\begin{equation*}
    D_N(t)=\int_{\mathcal Z}\Phi_t(z)\,\mathcal N_N(dz), \qquad 0\leq t\leq T.
\end{equation*}

We next verify the square-integrability of the envelope $G_T$. 
% Recall that, for a conversation mark $\zeta$, we let $L(\zeta):=\sum_{\ell=1}^{M-1}W_\ell$ denote the elapsed time between its first and last prompts. By definition,
% \begin{equation*}
%     G_T(a,\zeta)=B(\zeta)\one\{-L(\zeta)<a\leq T\}.
% \end{equation*}
For each fixed mark $\zeta$, the set of arrival times $a$ for which $G_T(a,\zeta)$ is nonzero has Lebesgue measure $T+L(\zeta)$. Therefore,
\begin{align*}
    V_T &=\int_{\mathcal Z}G_T(z)^2\,\nu(dz)\\
    &=\lambda_0\E\left[B^2(T+L)\right]\\
    &=\lambda_0T\E[B^2]+\lambda_0\mu_W\E[(M-1)B^2],
\end{align*}
where we obtain the last equality because the interarrival times are i.i.d. with mean $\mu_W$ and are independent of $(M,\{X_\ell,Y_\ell\}_{\ell\geq1})$. In addition, by \cref{ass: Regularity of one cumulative conversation blocks}, the finiteness of $B$'s M.G.F. around $0$ implies that all polynomial moments of $B$ are finite. Moreover, because every turn contributes at least one KV block, we have $M\leq B$ almost surely. Hence, 
\begin{equation*}
    \E[(M-1)B^2]\leq\E[B^3]<\infty,
\end{equation*}
and we conclude \eqref{equ: PRM second moment constant} because
\begin{equation*}
    V_T=\lambda_0T\E[B^2]+\lambda_0\mu_W\E[(M-1)B^2]<\infty.
\end{equation*}

Finally, because the path functional $\Phi_t(z)$ depends only on the background conversation trajectory and not on the tagged turn state $j$, the resulting displacement process has the same law for every fixed $j$ with $S_j>0$. Hence,
\begin{equation*}
    D_{N,j}(\cdot)\stackrel{\mathcal L}{=}D_N(\cdot),
\end{equation*}
which completes the proof.
\end{proof}


\subsubsection{Displacement Process Decomposition and the Mean Displacement Curve}
\label{sssec: Displacement Process Decomposition}

Fix a turn index $j$ with $S_j>0$. Under the stationary Palm distribution, suppose that a tagged conversation is accessed at its turn $j$ at time $0$ and moved to the MRU end. Recall that $D_{N,j}(t)$ denotes the corresponding displacement process and that, by \cref{lem: Poisson Random Measure Representation},
\begin{equation*}
    D_{N,j}(\cdot)\stackrel{\mathcal L}{=}D_N(\cdot),
\end{equation*}
where the law of $D_N(\cdot)$ is independent of the tagged turn state $j$.

We now decompose the representative displacement process $D_N(t)$ into two subprocesses, similarly to the proof of \cref{lem: Poisson Random Measure Representation}:
\begin{equation}
    \label{equ: displacement decomposition}
    D_N(t)=D_N^{\init}(t)+D_N^{\newb}(t).
\end{equation}
Here, $D_N^{\init}(t)$ denotes the displacement contributed by pre-existing history blocks of background conversations that are already active at time $0$ and are moved ahead of the tagged content during the time period $(0,t]$, during which their subsequent access happens. Additionally, $D_N^{\newb}(t)$ denotes the displacement contributed by newly generated KV blocks during $(0,t]$.


We first characterize the mean curve of the normalized pre-existing history contribution $\frac{D_N^{\init}(t)}{N}$.

\begin{lemma}[Mean Normalized Initial-Prefix Contribution]
    \label{lem: Initial-prefix contribution}
    During the time period $(0,t]$, the total pre-existing history content mass contributed to the displacement process, $D_N^{\init}(t)$, has the following normalized mean
    \begin{equation}
    \label{equ: mean normalized initial-prefix contribution}
        \frac{1}{N}\E[D_N^{\init}(t)]=\lambda_0A_H\int_0^t\bar F(s)\,ds,
    \end{equation}
    where $\bar F(t):=1-F(t)$ denotes the tail probability of the prompt interarrival distribution and $A_H:=\sum_{j\geq1}S_{j+1}\bar R_j$ denotes the expected amount of reusable history content, as defined in \eqref{equ: number of blocks eligible for reuse}.
\end{lemma}


\begin{proof}
    Fix a turn state $\ell\geq1$. A conversation that reaches turn $\ell+1$ has a waiting interval of length $W_\ell$ after completing turn $\ell$. 
    
    Let $s<0$ denote the time at which this conversation completes turn $\ell$, and let its inter-turn waiting time be $W_\ell=w$. Then its waiting interval is $[s,s+w)$. For its pre-existing history content to contribute to the displacement during $(0,t]$, the conversation's next prompt must arrive no later than time $t$. Equivalently, we require
    \begin{equation*}
        s<0<s+w\leq t.
    \end{equation*}
    For a fixed $w$, the set of waiting times $s$ satisfying this condition is
    \begin{equation*}
        -w<s\leq\min\{0,t-w\},
    \end{equation*}
    the admissible length of which is $\min\{t,w\}$. Therefore, when we integrate over the stationary time shifts of the marked conversation trajectories, the contribution associated with an inter-turn interval of length $w$ is weighted by $\min\{t,w\}$. 
    
    To determine how frequently such waiting intervals occur on the stationary time axis, we recall that conversations arrive according to a Poisson process with rate $N\lambda_0$. Among these conversations, a fraction $S_{\ell+1}=\pro(M\geq \ell+1)$ reaches turn $\ell+1$. In addition, the push-forward time $\sum\limits_{j=1}^{\ell} W_j$ for each conversation is i.i.d.. Therefore, by the mapping theorem of Poisson process, the arrival process of active turn-$\ell$ conversations that can reach turn $\ell+1$ is a Poisson process with rate $N\lambda_0 S_\ell p_{\ell} = N\lambda_0 S_{\ell+1}$.
    
    % By Campbell's theorem, the expected number of such waiting intervals that contain time $0$ and terminate during $(0,t]$ is
    % \begin{align*}
    %     N\lambda_0S_{\ell+1}\E\left[\min\{t,W_\ell\}\right]
    %     &=N\lambda_0S_{\ell+1}\int_0^t\pro(W_\ell>u)\,du\\
    %     &=N\lambda_0S_{\ell+1}\int_0^t\bar F(u)\,du.
    % \end{align*}
    
    
    % Consequently, the expected number of turn-$\ell$ waiting intervals whose starting points fall in an infinitesimal interval $ds$ is
    % \begin{equation*}
    %     N\lambda_0S_{\ell+1}\,ds.
    % \end{equation*}
    
    Therefore, by Campbell's theorem, we obtain the expected number of turn-$\ell$ waiting intervals that are active at time $0$ and terminate during $(0,t]$:
    \begin{align*}
        \E \bigg[\int_{-W_\ell}^{\min\{0,t-W_\ell\}} N\lambda_0 S_{\ell+1} ds \bigg ] &=  N\lambda_0S_{\ell+1}\E[\min\{t,W_\ell\}]\\
        & =N\lambda_0S_{\ell+1}\int_0^t\pro(W_\ell>u)\,du \\&=N\lambda_0S_{\ell+1}\int_0^t\bar F(u)\,du.
    \end{align*}

    At the terminating prompt of each such interval, its pre-existing history content of size $R_\ell$ is moved in front of the tagged content. Because $W_\ell$ is independent of the conversation length and block increments, the expected turn-$\ell$ history contribution is therefore
    \begin{equation*}
        N\lambda_0S_{\ell+1}\bar R_\ell\int_0^t\bar F(u)\,du.
    \end{equation*}
    Summing over all  $\ell\geq1$ and recalling that $A_H=\sum_{\ell\geq1}S_{\ell+1}\bar R_\ell$, we obtain that
    \begin{equation*}
        \E[D_N^{\mathrm{init}}(t)] = N\lambda_0A_H\int_0^t\bar F(u)\,du.
    \end{equation*}
\end{proof}


Similarly, the normalized new blocks contribution $\frac{D_N^{\newb}(t)}{N}$ admits the following expectation, where we recall from \eqref{equL The history content size R_j} and \eqref{equ: cumulative number of blocks} that $B: = R_M$ denotes the total lifetime KV block numbers of a conversation.

\begin{lemma}[Mean Normalized New Blocks Contribution]
\label{lem: Mean New Blocks Contribution}
    For every $t\geq 0$, we let $D_N^{\newb}(t)$ denote the new blocks generated by new prompts in the time duration $(0,t]$. Then the mean of this normalized contribution is 
    \begin{equation}
        \label{equ: Mean New Blocks Contribution}
        \frac{1}{N}\E[D_N^{\newb}(t)] = \lambda_0\E[B]t.
    \end{equation}
\end{lemma}

\begin{proof}
    In stationarity, an active conversation at turn state $j$ generates a new prompt, with mean number of blocks $\E [Z_{j+1}] = \E[X_{j+1}+Y_{j}]$. In addition, recall our argument in the proof of \cref{lem: Initial-prefix contribution}, that the arrival of these turn-$(j+1)$ prompts forms a Poisson process with rate $N \lambda_0 S_{j+1}$. In addition, the new conversation has arrival rate $N \lambda_0$, each contributes blocks with mean $\E [Z_1] = \E X_1$. Therefore, the total new block growth rate due to new block arrivals is 
    \begin{equation}
    \label{equ: new block growth rate}
        \sum\limits_{j\geq 1} N\lambda_0 S_{j} \E[Z_j] =  \sum\limits_{j\geq 1} N\lambda_0S_{j+1}\E [X_{j+1}+Y_{j}] + N \lambda_0 \E [X_1]
    \end{equation}
    
    Because $B=R_M$, by \cref{ass:independent block incrememts}, we have that 
    \begin{align*}
        \E[B] & = \E[X_1] + \E\Big[ \sum\limits_{j=1}^{M-1} X_{j+1}+Y_j \Big]\\
        &  = \E [X_1] + \sum_{j\geq1} \E[ (X_{j+1}+Y_j)\one\{M\geq j+1\} ] \\
        & = \E [X_1] + \sum_{j\geq1} S_{j+1}\E[X_{j+1}+Y_{j}].
    \end{align*}
    Therefore, we conclude \eqref{equ: Mean New Blocks Contribution}.
\end{proof}


Now we combine \cref{lem: Initial-prefix contribution,lem: Mean New Blocks Contribution} to obtain the following mean displacement curve.

\begin{lemma}[Mean Displacement Curve]
\label{lem: Mean Displacement Curve}
    For every $t\geq0$, the displacement process $D_{N,j}(t)$ has the same law as $D_N(t)$ for every fixed tagged turn state $j$ with $S_j>0$. Moreover,
    \begin{equation}
    \label{equ: Mean Displacement Curve}
        \frac{1}{N}\E[D_N(t)]=d(t):= \lambda_0\E[B]t + \lambda_0A_H\int_0^t\bar F(s)\,ds,
    \end{equation}
    where $\bar F(t):=1-F(t)$ denotes the tail probability of $F$. Moreover, $d$ is strictly increasing and Lipschitz-continuous, with
    \begin{equation*}
        \lambda_0\E[B](t-s) \leq d(t)-d(s) \leq \lambda_0(\E[B]+A_H)(t-s), \qquad\text{for all } 0\leq s\leq t.
    \end{equation*}
\end{lemma}

\begin{proof}
    We obtain \eqref{equ: Mean Displacement Curve} by directly combining \cref{lem: Initial-prefix contribution,lem: Mean New Blocks Contribution}. 

    Because $0\leq\bar F(u)\leq1$, for every $0\leq s\leq t$,
    \begin{align*}
        d(t)-d(s) =\lambda_0\E[B](t-s)+ \lambda_0A_H\int_s^t\bar F(u)\,du,
    \end{align*}
    and hence
    \begin{equation*}
        \lambda_0\E[B](t-s) \leq d(t)-d(s) \leq \lambda_0(\E[B]+A_H)(t-s).
    \end{equation*}
    Thus $d$ is strictly increasing and Lipschitz-continuous.   
\end{proof}

\subsubsection{Sample Path Convergence of the Normalized Displacement Process}
\label{sssec: normalized displacement sample path convergence}

We are now prepared to prove \cref{thm: Displacement Process Convergence}, the finite-size concentration and sample path convergence of the normalized displacement process.

By \cref{lem: Poisson Random Measure Representation}, for every fixed turn state $j$ with $S_j>0$, the process $\{D_{N,j}(t):0\leq t\leq T\}$ has the same law as the canonical displacement process $\{D_N(t):0\leq t\leq T\}$. We therefore omit the tagged turn-state index $j$ and study $D_N(t)$ directly.


\begin{reptheorem}{thm: Displacement Process Convergence}[Sample-Path Convergence of the Normalized Displacement Process]
    The tagged turn-$j$ content process follows the same law, such that $D_{N,j}(\cdot)\overset{d}{=}D_N(\cdot)$ for every turn state $j$. In addition, for a fixed $T<\infty$, let $V_T  = \lambda_0T\E[B^2]+\lambda_0\mu_W\E[(M-1)B^2]<\infty.$, and let $L :=\sup_{0\leq t\leq T}d'(t)=\lambda_0\E[B]+\lambda_0A_H$.
    Then, for every $t\in[0,T]$ and every $\varepsilon>0$,
    \begin{equation*}
        \pro\left(\left|\frac{D_N(t)}{N}-d(t)\right|>\varepsilon\right)\leq\frac{V_T}{N\varepsilon^2}.
    \end{equation*}
    Moreover, for every $0<\varepsilon\leq1$,
    \begin{equation*}
        \pro\left(\sup_{0\leq t\leq T}\left|\frac{D_N(t)}{N}-d(t)\right|>\varepsilon\right)\leq\frac{C_T}{N\varepsilon^3},
    \end{equation*}
    where $C_T:=\frac{16V_T}{9}(2+4 L T).$
    
    Consequently, for any fixed $T\geq 0$,
    \begin{equation*}
        \sup_{0\leq t\leq T}\left|\frac{D_N(t)}{N}-d(t)\right|\xrightarrow{\mathbb P}0.
    \end{equation*}
    In particular, the finite-size bounds imply
    \begin{equation*}
        \left|\frac{D_N(t)}{N}-d(t)\right|=O_{\mathbb P}(N^{-1/2})
    \end{equation*}
    for every fixed $t$, and
    \begin{equation*}
        \sup_{0\leq t\leq T}\left|\frac{D_N(t)}{N}-d(t)\right|=O_{\mathbb P}(N^{-1/3}).
    \end{equation*}
\end{reptheorem}
\begin{proof}

    For every $t\in[0,T]$, the Poisson random measure representation in \cref{lem: Poisson Random Measure Representation} gives that
    \begin{equation}
    \label{equ: displacement variance exact}
        \text{Var}\left(\frac{D_N(t)}{N}\right)=\frac{1}{N}\int_{\mathcal Z}\Phi_t^2(z)\,\nu(dz) \leq \frac{1}{N}\int_{\mathcal Z}G_T^2(z)\,\nu(dz) =\frac{V_T}{N}.
    \end{equation}
    In addition, due to \cref{lem: Mean Displacement Curve}, $\frac{1}{N}\E[D_N(t)] = d(t)$, we have that 
    \begin{equation}
        \label{equ: displacement MSE upper bound}
        \sup_{0\leq t\leq T}\E\left[\left|\frac{D_N(t)}{N}-d(t)\right|^2\right]\leq\frac{V_T}{N}.
    \end{equation}
    Therefore, applying Chebyshev's inequality, we first derive the following pointwise second-moment bound for any $\varepsilon>0$:
    \begin{equation*}
        \pro\left(\left|\frac{D_N(t)}{N}-d(t)\right|>\varepsilon\right)\leq\frac{1}{\varepsilon^2}\E\left[\left|\frac{D_N(t)}{N}-d(t)\right|^2\right]\leq\frac{V_T}{N\varepsilon^2}.
    \end{equation*}

    We next demonstrate the upper bound for the sample-path $\mathbb L^\infty$-difference. Recall from \cref{lem: Mean Displacement Curve} that for any $t>0$,
    \begin{equation*}
        d'(t)=\lambda_0\E[B]+\lambda_0A_H \bar F(t) \leq L := d'(0)=\lambda_0\E[B]+\lambda_0A_H.
    \end{equation*}
    Hence, the mean displacement curve $d$ is $L$-Lipschitz.

    Fix $0<\varepsilon\leq1$. Choose an equal partition
    \begin{equation*}
        \Pi = \{t_0,\cdots,t_K\},\quad \text{where}\quad 0=t_0<t_1<\cdots<t_K=T,
    \end{equation*}
    satisfying that
    \begin{equation*}
    \|\Pi\| : =\max_{0\leq k<K}(t_{k+1}-t_k)\leq\frac{\varepsilon}{4 L }.
    \end{equation*}
    It is enough to choose $K+1\leq 2+\frac{4 L T}{\varepsilon}$, i.e., $K=  \lfloor1+\frac{4 L T}{\varepsilon}\rfloor$.

    Fix $t\in[t_k,t_{k+1}]$. Because $D_N(t)$ and  $d(t)$ are both nondecreasing,
    \begin{equation*}
        D_N(t_k)\leq D_N(t)\leq D_N(t_{k+1}) \quad \text{and} \quad d(t_k)\leq d(t)\leq d(t_{k+1}).
    \end{equation*}
    Therefore, 
    \begin{align*}
        \frac{D_N(t)}{N}-d(t)
        &\leq\frac{D_N(t_{k+1})}{N}-d(t_k) \\
        &=\left(\frac{D_N(t_{k+1})}{N}-d(t_{k+1})\right)+\left(d(t_{k+1})-d(t_k)\right) \\
        &\leq\left|\frac{D_N(t_{k+1})}{N}-d(t_{k+1})\right|+\frac{\varepsilon}{4}.
    \end{align*}
    and similarly,
    \[
    d(t)-\frac{D_N(t)}{N} \leq \left|\frac{D_N(t_{k})}{N}-d(t_{k})\right|+\frac{\varepsilon}{4}.
    \]

    To establish 
    \begin{equation*}
        \sup_{0\leq t\leq T}\left|\frac{D_N(t)}{N}-d(t)\right|\leq\varepsilon,
    \end{equation*}
    it is enough to let, for any $1\leq k\leq K$ 
    \begin{equation*}
        \max_{0\leq k\leq K}\left|\frac{D_N(t_k)}{N}-d(t_k)\right|\leq\frac{3\varepsilon}{4}.
    \end{equation*}
    Therefore, we can bound the $L^\infty$-distance via a union bound, such that
    \begin{align*}
        \pro\left(\sup_{0\leq t\leq T}\left|\frac{D_N(t)}{N}-d(t)\right|>\varepsilon\right)& \leq\sum_{k=0}^{K}\pro\left(\left|\frac{D_N(t_k)}{N}-d(t_k)\right|>\frac{3\varepsilon}{4}\right)\\
        & \leq (K+1) \frac{16V_T}{9N\varepsilon^2}\\
        & \leq (2+\frac{4 L T}{\varepsilon})\frac{16V_T}{9N\varepsilon^2}.
    \end{align*}

    When $0<\varepsilon\leq1$, the last term can be further bounded by  
    \begin{equation*}
        (2+\frac{4 L T}{\varepsilon})\frac{16V_T}{9N\varepsilon^2} \leq \frac{2+4 L T}{\varepsilon}\frac{16V_T}{9N\varepsilon^2} = \frac{C_T}{N\varepsilon^3}.
    \end{equation*}  


    We finally show the convergence result with the given rate. For the pointwise convergence rate, set $\varepsilon = \delta/\sqrt N$ for some fixed $\delta>0$ in \eqref{equ: displacement pointwise concentration}. Then
    \begin{equation*}
        \pro\left(\sqrt N\left|\frac{D_N(t)}{N}-d(t)\right|>\delta \right)\leq\frac{V_T}{\delta^2}.
    \end{equation*}
    Taking $\limsup_{N\to\infty}$ and then $\delta\to\infty$, we have that
    \begin{equation*}
        \left|\frac{D_N(t)}{N}-d(t)\right|=O_{\mathbb P}(N^{-1/2}).
    \end{equation*}
    
    For the sample-path convergence rate, set $\varepsilon=\delta N^{-1/3}$ in \eqref{equ: displacement sample path concentration}. For every fixed $\delta>0$ and all sufficiently large $N$, $\varepsilon\leq 1$, and
    \begin{equation*}
        \pro\left(N^{1/3}\sup_{0\leq t\leq T}\left|\frac{D_N(t)}{N}-d(t)\right|>\delta\right)\leq\frac{C_T}{\delta^3}.
    \end{equation*}
    Taking $\limsup_{N\to\infty}$ and then $\delta\to\infty$, we have that
    \begin{equation*}
        \sup_{0\leq t\leq T}\left|\frac{D_N(t)}{N}-d(t)\right|=O_{\mathbb P}(N^{-1/3}).
    \end{equation*}

    Finally, because $D_{N,j}(\cdot)\stackrel{\mathcal L}{=}D_N(\cdot)$ for every fixed turn state $j$ with $S_j>0$, the same concentration bounds, convergence in probability, and finite-size rates hold for the normalized tagged-conversation displacement process $\frac{D_{N,j}(\cdot)}{N}$. This completes the proof.
\end{proof}

\subsection{Proof of \cref{thm: Characteristic Time Convergence- Full}}
\label{ssec: Proof of the Full Version of CT convergence}

We next focus on proving the convergence of the eviction time of a tagged turn-$j$ conversation content. First, we bound the tail of the turn-$j$ conversation content length $R_j$ for each $j\geq 1$, via the following lemma.


\begin{lemma}[Tail Bound of Turn-$j$ Content Length]
\label{lem: Tail Bound of Turn-j Content}
    Fix a turn state $j$ such that $S_j>0$. For any fixed $\theta\in(0,\theta_0)$, we have the following tail bound
    \begin{equation}
    \label{equ: Rj exponential tail}
        \pro(R_j>x)\leq \frac{\E[e^{\theta B}]}{S_j}e^{-\theta x}, \qquad x\geq0.
    \end{equation}
    Consequently,
    \begin{equation}
    \label{equ: Rj order}
        R_j=O_{\mathbb P}(1).
    \end{equation}
\end{lemma}

\begin{proof}

    By \cref{ass:independent block incrememts}, the block increments are independent of $M$. Because $R_j$ depends only on the block increments up to turn $j$, $R_j$ is independent of $M$. Hence,
    \begin{equation*}
        \E[e^{\theta R_j}]=\E[e^{\theta R_j}\mid M\geq j].
    \end{equation*}
    Conditional on the event $\{M\geq j\}$, the turn-$j$ history content is contained in the total lifetime block mass, and therefore $R_j\leq B$. Hence,
    \begin{align*}
        \E[e^{\theta R_j}]
        =\frac{1}{S_j}\E\left[e^{\theta R_j}\one\{M\geq j\}\right] 
        \leq\frac{1}{S_j}\E\left[e^{\theta B}\one\{M\geq j\}\right] \leq\frac{\E[e^{\theta B}]}{S_j}.
    \end{align*}
    Applying Markov's inequality, we obtain the following bound for every $x>0$:
    \begin{equation*}
        \pro(R_j>x)\leq e^{-\theta x}\E[e^{\theta R_j}]\leq\frac{\E[e^{\theta B}]}{S_j}e^{-\theta x}.
    \end{equation*}
     Because the distribution of $R_j$ is independent of $N$, the family $\{R_j\}_{N\geq1}$ is tight, and hence $R_j=O_{\mathbb P}(1)$.
\end{proof}


We are now ready to prove \cref{thm: Characteristic Time Convergence- Full}.

\begin{reptheorem}{thm: Characteristic Time Convergence- Full}[Characteristic Time Convergence]
    Fix a turn state $j$ such that $S_j>0$, and a finite horizon $T>T_C$. For every
    \begin{equation*}
        0<\varepsilon<\min\{T_C,T-T_C\},
    \end{equation*}
    the characteristic time satisfies
    \begin{equation*}
        \pro\left(\left|\tau_C^{(N)}(j)-T_C\right|>\varepsilon\right)
        \leq\frac{18V_T}{Nc_d^2\varepsilon^2}
        +\frac{\E[e^{\theta B}]}{S_j}\exp\left(-\frac{\theta c_d}{3}N\varepsilon\right) = p_{N,j}(\varepsilon),
    \end{equation*}
    where $c_d:=\lambda_0\E[B]$ and $V_T:=\int_{\mathcal Z}G_T(z)^2\,\nu(dz)<\infty$ is the second-moment constant introduced in \cref{lem: Poisson Random Measure Representation}. Consequently,
    \begin{equation*}
        \tau_C^{(N)}(j)\xrightarrow{\mathbb P}T_C
    \end{equation*}
    for every fixed turn state $j$, and
    \begin{equation}
    \label{equ: characteristic time convergence rate-2}
        \tau_C^{(N)}(j)-T_C=O_{\mathbb P}(N^{-1/2}).
    \end{equation}
\end{reptheorem}

\begin{proof}
    First, because for any $t>0$, $d'(t) = c_d + \lambda_0 A_H \bar F(t) \geq c_d$, thus for every $0<\varepsilon<T_C$,
    \begin{equation}
    \label{equ: lower separation around TC}
        d(T_C-\varepsilon)\leq1-c_d\varepsilon,
    \end{equation}
    and for every $\varepsilon>0$,
    \begin{equation}
    \label{equ: upper separation around TC}
        d(T_C+\varepsilon)\geq1+c_d\varepsilon.
    \end{equation}

    Fix $0<\varepsilon<\min\{T_C,T-T_C\}$ and define
    \begin{equation*}
        t_-:=T_C-\varepsilon, \qquad t_+:=T_C+\varepsilon, \qquad \eta:=\frac{c_d\varepsilon}{3}.
    \end{equation*}
    Then we rewrite equations \eqref{equ: lower separation around TC} and \eqref{equ: upper separation around TC}, such that 
    \begin{equation*}
    d(t_-)\leq1-3\eta, \quad \text{and }\quad d(t_+)\geq1+3\eta, \quad \text{respectively.}
    \end{equation*}
    
    We introduce the event
    \begin{equation*}
        \mathcal E_N(\varepsilon) := \left\{\left|\frac{D_N(t_-)}{N}-d(t_-)\right |\leq\eta\right\}\cap\left\{\left|\frac{D_N(t_+)}{N}-d(t_+)\right|\leq\eta\right\}\cap\left\{\frac{R_j}{N}\leq\eta\right\},
    \end{equation*}
    and demonstrate that on $\mathcal E_N(\varepsilon)$, we have that 
    \begin{equation}
        \label{equ: Bound the tauNCj}
        T_C-\varepsilon< \tau_C^{(N)}(j)\leq T_C+\varepsilon.
    \end{equation}
    
    On the one hand, when event $\mathcal E_N(\varepsilon)$ happens,
    \begin{equation*}
        \frac{D_N(t_-)}{N}\leq d(t_-)+\eta\leq1-2\eta, \quad \text{and} \quad  \frac{N-R_j}{N}=1-\frac{R_j}{N}\geq1-\eta,
    \end{equation*}
    which implies that, at the time $t_-$, the displacement is below the threshold, i.e., $D_N(t_-)<N-R_j$, and thus $\tau_C^{(N)}(j)>T_C-\varepsilon.$

    On the other hand, at time $t_+$, given the event $\mathcal E_N(\varepsilon)$ happening, 
    \begin{equation*}
        \frac{D_N(t_+)}{N}\geq d(t_+)-\eta\geq1+2\eta>1.
    \end{equation*}
    Because $R_j > 0$, $N-R_j< N<D_N(t_+)$, and hence $\tau_C^{(N)}(j)\leq T_C+\varepsilon.$ Thus we conclude the inclusion $\mathcal E_N(\varepsilon)\subseteq\left\{\left|\tau_C^{(N)}(j)-T_C\right|\leq\varepsilon\right\}$ and claim \eqref{equ: Bound the tauNCj}.

    We now apply the union bound, which gives
    \begin{equation}
        \label{equ: union bound for tau and TC diff}
        \pro\left(\left|\tau_C^{(N)}(j)-T_C\right|>\varepsilon\right)
        \leq\pro\left(\left|\frac{D_N(t_-)}{N}-d(t_-)\right|>\eta\right)+\pro\left(\left|\frac{D_N(t_+)}{N}-d(t_+)\right|>\eta\right)+\pro(R_j>N\eta).
    \end{equation}
    
    We use the pointwise displacement tail bound in \cref{thm: Displacement Process Convergence} to bound the first two terms in the R.H.S. of \eqref{equ: union bound for tau and TC diff}, such that
    \begin{equation*}
        \pro\left(\left|\frac{D_N(t_\pm)}{N}-d(t_\pm)\right|>\eta\right)\leq\frac{V_T}{N\eta^2} \stackrel{\eta=c_d\varepsilon/3}{=}\frac{9V_T}{N c_d^2\varepsilon^2}.
    \end{equation*}
    In addition, the last term on the R.H.S. of \eqref{equ: union bound for tau and TC diff} can be bounded using \cref{lem: Tail Bound of Turn-j Content}:
    \begin{equation*}
        \pro(R_j>N\eta)\leq\frac{\E[e^{\theta B}]}{S_j}e^{-\theta N\eta}=\frac{\E[e^{\theta B}]}{S_j}\exp\left(-\frac{\theta c_d}{3}N\varepsilon\right).
    \end{equation*}
    Combining these bounds proves \eqref{equ: characteristic time finite N bound}. Moreover, for every fixed $\varepsilon>0$, all terms on the right-hand side of \eqref{equ: union bound for tau and TC diff} converge to zero as $N\rightarrow\infty$. Hence,
    \begin{equation*}
        \tau_C^{(N)}(j)\xrightarrow{\mathbb P}T_C.
    \end{equation*}
    
    Finally, the convergence rate can be shown as follows. When we let $x>0$ be fixed and set $\varepsilon=\frac{x}{\sqrt N}$, $\varepsilon<\min\{T_C,T-T_C\}$ can always be satisfied for sufficiently large $N$, and \eqref{equ: characteristic time finite N bound} gives
    \begin{equation*}
        \pro\left(\sqrt N\left|\tau_C^{(N)}(j)-T_C\right|>x\right)\leq\frac{18V_T}{c_d^2x^2}+\frac{\E[e^{\theta B}]}{S_j}\exp\left(-\frac{\theta c_d}{3}x\sqrt N\right).
    \end{equation*}
    Therefore, letting $N \uparrow \infty$, we conclude that $\tau_C^{(N)}(j)-T_C=O_{\mathbb P}(N^{-1/2}).$
\end{proof}



\subsection{Proof of \cref{thm: Hit Ratio Convergence Full}}
\label{ssec: Proof of the Full Version of hit ratio conv MCM-LRU short}

We finally establish the mean-field convergence of the  hit ratio under the whole-content LRU policy. For every $j\geq2$, recall that $H_{N,j}$ denotes the indicator that the history content of size $R_{j-1}$ remains resident in HBM when the turn-$j$ prompt arrives.

To analyze the limiting behavior of the hit ratio, we first define some auxiliary notations. 

Recall that $A_H=\sum_{j\geq1}S_{j+1}\bar R_j$ denotes the expected amount of reusable history content workload. We further define $A_H^{(J)}$ as a truncation of $A_H$ at level $J\geq 1$, i.e.,
\begin{equation}
\label{equ: truncated AH J}
    A_H^{(J)}:=\sum_{j=1}^{J}S_{j+1}\bar R_j.
\end{equation}
The tail workload is then defined such that
\begin{equation}
\label{equ: reusable workload tail}
    \mathcal T_J:=\sum_{j>J}S_{j+1}\bar R_j=A_H-A_H^{(J)}.
\end{equation}
In addition, we let 
\begin{equation}
\label{equ: truncated hit second moment constant}
    K_J:=\sum_{j=1}^{J}S_{j+1}\sqrt{\E[R_j^2]}
\end{equation}
which will be used to control the finite-$J$ approximation error.


We prove \cref{thm: Hit Ratio Convergence Full} through the following lemmas. We first establish an exact representation of the finite-$N$ LRU hit ratio in terms of the characteristic time. Recall that $W_j$, the prompt interarrival time between turns $j$ and $j+1$, has cumulative distribution function $F$. By the definition of the characteristic time, we have that
\begin{equation}
\label{equ: hit indicator characteristic time representation}
    H_{N,j+1}=\one\left\{W_j<\tau_C^{(N)}(j)\right\}.
\end{equation}


\begin{lemma}[Exact Representation of the LRU Hit Ratio]
\label{lem: Exact Representation of LRU Hit Probability}
For every $N$, the  hit ratio admits the representation
\begin{equation}
\label{equ: exact LRU hit ratio characteristic representation}
    h ^{(N)} = \frac{1}{A_R}\sum_{j\geq1}S_{j+1}\E\left[R_j F(\tau^{N}_C(j)) \right].
\end{equation}
\end{lemma}

\begin{proof}
    By \cref{ass: Exponential Prompt Inter-arrival Time}, the tagged-conversation interarrival time $W_j$ is independent of the block increments and hence of $R_j$. Moreover, the displacement process is determined entirely by the background conversations, which are independent of the tagged conversation. Therefore, $W_j$ is independent of $(R_j,\tau_C^{(N)}(j))$. Conditioning on $(R_j,\tau_C^{(N)}(j))$ gives
    \begin{align*}
        \E\left[H_{N,j+1}\mid R_j,\tau_C^{(N)}(j)\right]
        =\pro\left(W_j<\tau_C^{(N)}(j)\mid R_j,\tau_C^{(N)}(j)\right) = F \big (\tau^{(N)}_C(j)\big) %\\&=1-e^{-r\tau_C^{(N)}(j)}.
    \end{align*}
    
    Using \eqref{equ: finite N LRU hit ratio},
    \begin{equation*}
        \E\left[\sum_{j=1}^{M-1}R_jH_{N,j+1}\right]= \sum_{j\geq1}\E\left[R_jH_{N,j+1}\one\{M\geq j+1\}\right].
    \end{equation*}
    Additionally, according to \cref{ass:independent block incrememts}, $M$ is independent of the block increment. Hence,
    \begin{equation}
    \label{equ: conditional hit probability}
        \E\Big[R_jH_{N,j+1}\one\{M\geq j+1\}\Big] =S_{j+1}\E\left[R_jH_{N,j+1}\right] =S_{j+1}\E\left[R_jF(\tau^{N}_C(j))\right].
    \end{equation}
    By summing over $j\geq 1$ and dividing by $A_R$, we prove \eqref{equ: exact LRU hit ratio characteristic representation}.
\end{proof}

Recall in \cref{thm: Characteristic Time Convergence- Full}, that the convergence of $\tau^{(N)}_C(j)$ is not uniform in $j$. We therefore first use the finite-turn proxy
\begin{equation}
    \label{equ: finite turn proxy J}
    \sum_{j=1}^{J}S_{j+1}\E\left[R_jF(\tau^{N}_C(j))\right].
\end{equation}
The difference between the finite-size and limiting hit ratios can then be decomposed into two parts: the tail contributions beyond turn $J$, which are controlled by $\mathcal T_J$, and the truncated approximation error $\left| \sum_{j=1}^{J}S_{j+1}\E\left[R_jF(\tau^{N}_C(j))\right]- A_H^{(J)}F(T_C) \right|$, which can be bounded via the following lemma.


\begin{lemma}[Finite-Turn Hit Ratio Error]
    \label{lem: Finite Turn Hit Ratio Error}
        Fix a finite $J$ such that $S_{J+1}>0$. For every $0<\varepsilon<\min\{T_C,T-T_C\}$,
    \begin{equation}
    \label{equ: finite turn hit ratio error}
        \left| \sum_{j=1}^{J}S_{j+1}\E\left[R_jF(\tau^{N}_C(j))\right]- A_H^{(J)}F(T_C)\right|\leq L_F \varepsilon A_H^{(J)}+K_J\sqrt{p_{N,J}(\varepsilon)}.
    \end{equation}
\end{lemma}



\begin{proof}
    By the triangle inequality, the L.H.S. of \eqref{equ: finite turn hit ratio error} can be bounded such that
    \begin{align*}
        \left|\sum_{j=1}^{J}S_{j+1}\E\left[R_jF\left(\tau_C^{(N)}(j)\right)\right]-\sum_{j=1}^{J}S_{j+1}\bar R_jF(T_C)\right| \leq\sum_{j=1}^{J}S_{j+1}\E\left[R_j\left|F\left(\tau_C^{(N)}(j)\right)-F(T_C)\right|\right],
    \end{align*}

    Due to \cref{ass: Exponential Prompt Inter-arrival Time}, $F(\cdot)$ is $L_F$-Lipschitz, i.e.,
    \begin{equation*}
        |F(t)-F(s)|\leq L_F|t-s|, \qquad \text{for all } 0\leq t<s<\infty.
    \end{equation*}
    For every $j\leq J$, we define the concentration event $\mathcal G_{N,j}(\varepsilon):=\left\{\left|\tau_C^{(N)}(j)-T_C\right|\leq\varepsilon\right\}$. On $\mathcal G_{N,j}(\varepsilon)$, the Lipschitz condition implies that
    \begin{equation*}
        \left|F\left(\tau_C^{(N)}(j)\right)-F(T_C)\right|\leq L_F\varepsilon.
    \end{equation*}
    Because $0\leq F(t)\leq1$ for all $t\geq0$, on the complementary event $\mathcal G_{N,j}(\varepsilon)^c$ we still have $\left|F\left(\tau_C^{(N)}(j)\right)-F(T_C)\right|\leq1$. Therefore,
    \begin{equation*}
        \E\left[R_j\left|F\left(\tau_C^{(N)}(j)\right)-F(T_C)\right|\right]
        \leq L_F\varepsilon\E[R_j]+\E\left[R_j\one\left\{\mathcal G_{N,j}(\varepsilon)^c\right\}\right].
    \end{equation*}
    By the Cauchy-Schwarz inequality,
    \begin{equation*}
        \E\left[R_j\one\left\{\mathcal G_{N,j}(\varepsilon)^c\right\}\right]
        \leq\sqrt{\E[R_j^2]}\sqrt{\pro\left(\left|\tau_C^{(N)}(j)-T_C\right|>\varepsilon\right)}.
    \end{equation*}

    According to \cref{thm: Characteristic Time Convergence- Full},
    \begin{equation*}
        \pro\left(\left|\tau_C^{(N)}(j)-T_C\right|>\varepsilon\right)\leq p_{N,j}(\varepsilon).
    \end{equation*}
    Moreover, because $S_j$ is nonincreasing in $j$, we have $S_j\geq S_J$ for every $j\leq J$. Hence, the finite-size bound in \cref{thm: Characteristic Time Convergence- Full} implies $p_{N,j}(\varepsilon)\leq p_{N,J}(\varepsilon),$
    and therefore
    \begin{equation*}
        \pro\left(\left|\tau_C^{(N)}(j)-T_C\right|>\varepsilon\right)\leq p_{N,J}(\varepsilon).
    \end{equation*}
    It follows that for each $1\leq j\leq J$,
    \begin{equation*}
        \E\left[R_j\left|F\left(\tau_C^{(N)}(j)\right)-F(T_C)\right|\right]
        \leq L_F\varepsilon\bar R_j+\sqrt{\E[R_j^2]}\sqrt{p_{N,J}(\varepsilon)}.
    \end{equation*}
    Multiplying by $S_{j+1}$ and summing over $j=1,\ldots,J$, we obtain that
    \begin{equation*}
        \sum_{j=1}^{J}S_{j+1}\E\left[R_j\left|F\left(\tau_C^{(N)}(j)\right)-F(T_C)\right|\right]
        \leq L_F\varepsilon A_H^{(J)}+K_J\sqrt{p_{N,J}(\varepsilon)},
    \end{equation*}
    and we prove \eqref{equ: finite turn hit ratio error}.
\end{proof}


We next control the error of the tail workload $\mathcal T_J = \sum\limits_{j>J}S_{j+1}\bar R_j$, via the lemma below.

\begin{lemma}[Bound for the Tail Workload]
\label{lem: Tail Bound for Reusable Workload}
    For every integer $J\geq1$ and any fixed $\theta\in(0,\theta_0)$, the tail workload $\mathcal T_J$ has the following upper bound:
    \begin{equation}
    \label{equ: reusable workload tail exponential bound}
        \mathcal T_J\leq \E[B^2e^{\theta B}]e^{-\theta(J+2)}.
    \end{equation}
    Consequently,
    \begin{equation}
    \label{equ: reusable workload tail convergence}
        \mathcal T_J\rightarrow 0.
    \end{equation}
\end{lemma}

\begin{proof}x
    Using the independence between $M$ and the block increments $R_j$'s, we can rewrite $\mathcal T_J$, such that
    \begin{align*}
        \mathcal T_J = \sum_{j>J}S_{j+1}\bar R_j =\sum_{j>J}\E\left[R_j\one\{M\geq j+1\}\right] =\E\Big[\sum_{j=J+1}^{M-1}R_j\Big].
    \end{align*}
    Conditional on the event $\{M\geq J+2\}$, each partial history content satisfies $R_j\leq B$. Hence,
    \begin{equation*}
        \sum_{j=J+1}^{M-1}R_j \leq MB\one\{M\geq J+2\}.
    \end{equation*}
    Because each turn generates at least one additional block, $M\leq B$ almost surely. Therefore,
    \begin{equation*}
        \sum_{j=J+1}^{M-1}R_j \leq B^2\one\{B\geq J+2\}.
    \end{equation*}
    For any fixed $\theta\in(0,\theta_0)$, because $\one\{B\geq J+2\} \leq e^{\theta (B-J-2)}$,
    \begin{equation*}
        B^2\one\{B\geq J+2\} \leq B^2e^{\theta B}e^{-\theta(J+2)} \quad \Rightarrow \quad  \mathcal T_J\leq\E[B^2e^{\theta B}]e^{-\theta(J+2)},
    \end{equation*}
    and we conclude \cref{lem: Tail Bound for Reusable Workload}.
\end{proof}

We now demonstrate proof of \cref{thm: Hit Ratio Convergence Full}, by applying \cref{lem: Exact Representation of LRU Hit Probability,lem: Finite Turn Hit Ratio Error,lem: Tail Bound for Reusable Workload}.

\begin{reptheorem}{thm: Hit Ratio Convergence Full}[Hit Ratio Convergence in the MCM-LRU Model]
    Under the MCM assumptions, fix a finite $J$ such that $S_{J+1}>0$, and fix $T>T_C$. For every $0<\varepsilon<\min\{T_C,T-T_C\}$,
    \begin{equation*}
    \left|h_{\LRU}^{(N)} -h_{\LRU}^{(\infty)}  \right| \leq \frac{L_F\varepsilon A_H^{(J)}}{A_R} +\frac{K_J}{A_R}\sqrt{p_{N,J}(\varepsilon)} +\frac{2\E[B^2e^{\theta B}]e^{-\theta(J+2)}}{A_R},
    \end{equation*}
    where $K_J:=\sum_{j=1}^{J}S_{j+1}\sqrt{\E[R_j^2]}$.
    Consequently, the system-level hit ratio under the whole-content LRU policy converges when $N\rightarrow\infty$:
    \begin{equation*}
        h_{\LRU}^{(N)} \rightarrow h_{\LRU}^{(\infty)}  =\frac{A_H}{A_R} F(T_C).
    \end{equation*}
\end{reptheorem}

    
\begin{proof}[Proof of \cref{thm: Hit Ratio Convergence Full}]
    Recall \cref{lem: Exact Representation of LRU Hit Probability} that the finite-size hit ratio has the following representation
    \begin{equation*}
    h ^{(N)} = \frac{1}{A_R} \sum_{j\geq1} S_{j+1} \E\left[R_j F\left(\tau_C^{(N)}(j)\right)\right],
    \end{equation*}
    and the limiting hit ratio,
    \begin{equation*}
    h ^{(\infty)} = \frac{1}{A_R} \sum_{j\geq1} S_{j+1}\bar R_j F(T_C).
    \end{equation*}
    We fix a truncation $J$ and use triangle inequality to bound the difference between $h ^{(N)}$ and $h ^{(\infty)}$, such that
    \begin{align*}
        \left|h ^{(N)} -h ^{(\infty)}  \right|
        &\leq \frac{1}{A_R}\left|\sum_{j=1}^{J}S_{j+1}\E\left[R_j F\left(\tau_C^{(N)}(j)\right)\right]-
        \sum_{j=1}^{J}S_{j+1}\bar R_j F(T_C)\right| \\
        &\quad+\frac{1}{A_R}\sum_{j>J}S_{j+1}\E\left[R_j F\left(\tau_C^{(N)}(j)\right)\right]+\frac{1}{A_R} \sum_{j>J}S_{j+1}\bar R_j F(T_C).
    \end{align*}
    Because $0\leq F(t)\leq1$,
    \begin{equation*}
        \sum_{j>J}S_{j+1}\E\left[R_j F\left(\tau_C^{(N)}(j)\right)\right]\leq \sum_{j>J}S_{j+1}\E\left[R_j\right] \leq \mathcal T_J,\quad\text{and} \quad \sum_{ j>J}S_{j+1}\bar R_j F(T_C)\leq\mathcal T_J.
    \end{equation*}
    Applying \cref{lem: Finite Turn Hit Ratio Error} to the finite-turn error, we obtain that
    \begin{equation*}
        \left|h ^{(N)} -h ^{(\infty)}  \right| \leq
        \frac{L_F \varepsilon A_H^{(J)}}{A_R} + \frac{K_J}{A_R}\sqrt{p_{N,J}(\varepsilon)} + \frac{2\mathcal T_J}{A_R}.
    \end{equation*}
    By \cref{lem: Tail Bound for Reusable Workload}, we have that $\mathcal T_J\leq\E[B^2e^{\theta B}]e^{-\theta(J+2)},$ which proves the upper bound \eqref{equ: finite N hit ratio error bound}.
    
    To prove the convergence result, we fix $J$ and choose $\varepsilon=N^{-1/4}$. For sufficiently large $N$, $\varepsilon<\min\{T_C,T-T_C\}$. The first term satisfies
    \begin{equation*}
        \frac{L_F \varepsilon A_H^{(J)}}{A_R}=O(N^{-1/4}).
    \end{equation*}
    Moreover,
    \begin{align*}
        p_{N,J}(N^{-1/4}) &= \frac{18V_T}{c_d^2N^{1/2}} + \frac{\E[e^{\theta B}]}{S_J} \exp\left(-\frac{\theta c_d}{3}N^{3/4}\right) = O(N^{-\frac{1}{2}}), 
    \end{align*}
    and therefore $\sqrt{p_{N,J}(N^{-1/4})}\rightarrow0.$
    
    Hence, for every fixed $J$,
    \begin{equation*}
        \limsup_{N\rightarrow\infty} \left|h ^{(N)} -h ^{(\infty)}  \right| \leq \frac{2\mathcal T_J}{A_R}.
    \end{equation*}
    Finally, by \cref{lem: Tail Bound for Reusable Workload}, $\mathcal T_J\rightarrow 0$ as $J\rightarrow\infty$. Taking $J\rightarrow\infty$, we conclude \cref{thm: Hit Ratio Convergence Full}.
\end{proof}

\subsection{Proof of \cref{thm: Practical Hit Ratio Convergence Full}}
\label{ssec: Proof of the Full Version of practical results}

We first show that the displacement-process analysis remains valid under the practical modification. For a fixed $T<\infty$, we define
\begin{equation*}
    V_T^{\mathrm{prac}}:=\int_{\mathcal Z}G_T(z)^2\,\nu(dz)=\lambda_0T\E[B^2]+\lambda_0\mu_W\E[(M-1)B^2]<\infty.
\end{equation*}
and
\begin{equation*}
    L_{\mathrm{prac}}:=\sup_{0\leq t < \infty}d_{\mathrm{prac}}'(t) =\lambda_0 \E[B]+\lambda_0 A_H^{\mathrm{hash}}.
\end{equation*}


\begin{lemma}[Normalized Practical Displacement Process Convergence]
\label{lem: Normalized Practical Displacement Process Convergence}
In the practical MCM, for every fixed $t\in[0,T]$ and $\varepsilon>0$, the deviation of the normalized displacement process has the following tail bound:
\begin{equation}
\label{equ: practical displacement pointwise concentration}
    \pro\left(\left|\frac{D_N^{\mathrm{prac}}(t)}{N}-d_{\mathrm{prac}}(t)\right|>\varepsilon\right) \leq\frac{V_T^{\mathrm{prac}}}{N\varepsilon^2}.
\end{equation}
For any fixed $t\leq T$, we have the pointwise convergence rate
\begin{equation}
\label{equ: practical displacement pointwise rate}
    \left|\frac{D_N^{\mathrm{prac}}(t)}{N}-d_{\mathrm{prac}}(t)\right|=O_{\mathbb P}(N^{-1/2}),
\end{equation}
and the uniform convergence admits the following convergence rate
\begin{equation}
\label{equ: practical displacement uniform rate}
    \sup_{0\leq t\leq T}\left|\frac{D_N^{\mathrm{prac}}(t)}{N}-d_{\mathrm{prac}}(t)\right|=O_{\mathbb P}(N^{-1/3}).
\end{equation}
In consequence,
\begin{equation}
\label{equ: practical displacement sample path convergence}
    \sup_{0\leq t\leq T}\left|\frac{D_N^{\mathrm{prac}}(t)}{N}-d_{\mathrm{prac}}(t)\right|\xrightarrow{\mathbb P}0,
\end{equation}

\end{lemma}

\begin{proof}
The practical system has the same stationary turn-state population and the same Poisson random measure structure as in \cref{lem: Poisson Random Measure Representation}. The only difference is that, when an existing history content is first moved ahead of the tagged conversation, its contribution changes from $R_j$ to $R_j-1$.

Accordingly, the practical displacement process admits the following  representation regardless of the turn state of the tagged conversation 
\begin{equation*}
    D_N^{\mathrm{prac}}(t)=\int_{\mathcal Z}\Phi_t^{\mathrm{prac}}(z)\,\mathcal N_N(dz),
\end{equation*}
where the practical contribution satisfies
\begin{equation*}
    0\leq\Phi_t^{\mathrm{prac}}(z)\leq\Phi_t(z)\leq G_T(z), \qquad 0\leq t\leq T.
\end{equation*}
Thus, we apply the same square-integrability argument as in \cref{lem: Poisson Random Measure Representation}.

The new-block contribution is unchanged and has normalized mean $\lambda_0\E[B]t$. For the initial-prefix contribution, an active turn-$j$ conversation contributes $R_j-1$ rather than $R_j$ when its next prompt arrives. Therefore,
\begin{equation*}
    \frac{1}{N}\E[D_N^{\mathrm{prac}}(t)]
    =\lambda_0\E[B]t+\lambda_0A_H^{\mathrm{hash}}\int_0^t\bar F(u)\,du
    =d_{\mathrm{prac}}(t).
\end{equation*}

Using the envelope $\Phi_t^{\mathrm{prac}}(z)\leq G_T(z)$, the same Poisson-integral variance calculation as in \cref{thm: Displacement Process Convergence}, we have the following variance bound
\begin{equation*}
    \text{Var}\left(\frac{D_N^{\mathrm{prac}}(t)}{N}\right)
    \leq\frac{V_T^{\mathrm{prac}}}{N}.
\end{equation*}
Hence, we claim \eqref{equ: practical displacement pointwise concentration} using Chebyshev's inequality.

Finally, $D_N^{\mathrm{prac}}(t)$ and $d_{\mathrm{prac}}(t)$ are nondecreasing, while $d_{\mathrm{prac}}$ is $L_{\mathrm{prac}}$-Lipschitz. Therefore, the partition and union-bound argument in the proof of \cref{thm: Displacement Process Convergence} applies without modification, which proves \eqref{equ: practical displacement sample path convergence} and \eqref{equ: practical displacement uniform rate}. The pointwise rate \eqref{equ: practical displacement pointwise rate} follows directly from \eqref{equ: practical displacement pointwise concentration}.
\end{proof}


We next consider the characteristic time of a reusable turn-$j$ history content.

\begin{lemma}[Practical Characteristic Time of a Turn-$j$ Tagged Conversation]
\label{lem: Practical Characteristic Time Turn-j}
    Fix a turn state $j$ such that $S_j>0$, and let
    \begin{equation*}
        m_{\mathrm{prac}}:=\lambda_0\E[B].
    \end{equation*}
    For every fixed $T>T_C^{\mathrm{prac}}$ and $0<\varepsilon<\min\{T_C^{\mathrm{prac}},T-T_C^{\mathrm{prac}}\},$ we have that
    \begin{equation}
    \label{equ: practical CT finite N bound}
        \pro\left(\left|\tau_{C,\mathrm{prac}}^{(N)}(j)-T_C^{\mathrm{prac}}\right|>\varepsilon\right) \leq \frac{18V_T^{\mathrm{prac}}}{Nm_{\mathrm{prac}}^2\varepsilon^2} + \frac{\E[e^{\theta B}]}{S_j} \exp\left(-\frac{\theta m_{\mathrm{prac}}}{3}N\varepsilon\right).
    \end{equation}
    Consequently, we have the practical characteristic time convergence, such that
    \begin{equation}
    \label{equ: practical CT convergence-2}
        \tau_{C,\mathrm{prac}}^{(N)}(j)\xrightarrow{\mathbb P}T_C^{\mathrm{prac}},
    \end{equation}
    and the convergence rate is given by
    \begin{equation}
    \label{equ: practical CT rate}
        \tau_{C,\mathrm{prac}}^{(N)}(j)-T_C^{\mathrm{prac}}=O_{\mathbb P}(N^{-1/2}).
    \end{equation}
\end{lemma}

\begin{proof}
    The proof follows the same threshold-sandwich argument as \cref{thm: Characteristic Time Convergence- Full}. The reusable tagged content now has size $R_j-1$ instead of $R_j$, and $0\leq R_j-1 \leq R_j.$ Hence, \cref{lem: Tail Bound of Turn-j Content} continues to control its normalized size $\frac{R_j-1}{N}$.
    
    Moreover, because the increasing rate of $d_{\mathrm{prac}}(t)$ has the lower bound
    \begin{equation*}
        d_{\mathrm{prac}}'(t)=m_{\mathrm{prac}}+\lambda_0A_H^{\mathrm{hash}}\bar F(t)
    \geq m_{\mathrm{prac}},
    \end{equation*}
    \begin{equation*}
        d_{\mathrm{prac}}(T_C^{\mathrm{prac}}-\varepsilon)\leq1-m_{\mathrm{prac}}\varepsilon, \quad \text{and} \quad  d_{\mathrm{prac}}(T_C^{\mathrm{prac}}+\varepsilon)\geq1+m_{\mathrm{prac}}\varepsilon.
    \end{equation*}
    
    We then similarly set $\eta:=\frac{m_{\mathrm{prac}}\varepsilon}{3}$ and repeat the event construction in the proof of \cref{thm: Characteristic Time Convergence- Full}, with $D_N$ replaced by $D_N^{\mathrm{prac}}$, $d$ replaced by $d_{\mathrm{prac}}$, and $R_j$ replaced by $R_j-1$. This gives
    \begin{align*}
        \pro\left(\left|\tau_{C,\mathrm{prac}}^{(N)}(j)-T_C^{\mathrm{prac}}\right|>\varepsilon\right)
        &\leq\pro\left(\left|\frac{D_N^{\mathrm{prac}}(T_C^{\mathrm{prac}}-\varepsilon)}{N}-d_{\mathrm{prac}}(T_C^{\mathrm{prac}}-\varepsilon)\right|>\eta\right) \\
        &\quad+ \pro\left(\left|\frac{D_N^{\mathrm{prac}}(T_C^{\mathrm{prac}}+\varepsilon)}{N}-d_{\mathrm{prac}}(T_C^{\mathrm{prac}}+\varepsilon)\right|>\eta\right)+ \pro(\frac{R_j-1}{N}>\eta).
    \end{align*}
    Using the pointwise concentration bound in \cref{lem: Normalized Practical Displacement Process Convergence}, we can bound the first two terms, such that
    \begin{equation*}
        \pro\left(\left|\frac{D_N^{\mathrm{prac}}(T_C^{\mathrm{prac}}\pm\varepsilon)}{N}
        -d_{\mathrm{prac}}(T_C^{\mathrm{prac}}\pm\varepsilon)\right|>\eta\right)
        \leq\frac{V_T^{\mathrm{prac}}}{N\eta^2}.
    \end{equation*}
    Additionally, because $R_j-1\leq R_j$, \cref{lem: Tail Bound of Turn-j Content} gives that
    \begin{equation*}
        \pro(R_j-1>N\eta)
        \leq\frac{\E[e^{\theta B}]}{S_j}e^{-\theta N\eta}.
    \end{equation*}
    Substituting $\eta=m_{\mathrm{prac}}\varepsilon/3$, we conclude \eqref{equ: practical CT finite N bound}. The convergence and $O_{\mathbb P}(N^{-1/2})$ rate then follow exactly as in the final part of the proof of \cref{thm: Characteristic Time Convergence- Full}.
\end{proof}


We finally establish the practical hit-ratio convergence. For a finite truncation level $J\geq1$, define the truncated $A_{H,\mathrm{hash}}$ at turn level $J$ as 
\begin{equation}
\label{equ: practical truncated AH}
    A_{H,\mathrm{hash}}^{(J)} :=\sum_{j=1}^{J}S_{j+1}\E[R_j-1],
\end{equation}
and the truncated second moment constant at level $J$,
\begin{equation}
\label{equ: practical truncated KJ}
    K_{J,\mathrm{hash}}:=\sum_{j=1}^{J}S_{j+1}\sqrt{\E[(R_j-1)^2]}.
\end{equation}

Similarly to \cref{ssec: Proof of the Full Version of hit ratio conv MCM-LRU short}, define the practical tail workload as
\begin{equation}
\label{equ: practical tail workload}
    \mathcal T_{J,\mathrm{hash}} :=\sum_{j>J}S_{j+1}\E[R_j-1].
\end{equation}
Because $0\leq R_j-1\leq R_j$, by \cref{lem: Tail Bound for Reusable Workload}, 
\begin{equation}
\label{equ: practical tail workload bound}
    0\leq\mathcal T_{J,\mathrm{hash}}\leq\mathcal T_J \leq\E[B^2e^{\theta B}]e^{-\theta(J+2)}.
\end{equation}

For fixed $T>T_C^{\mathrm{prac}}$, define the corresponding tail probability bound 
\begin{equation}
\label{equ: practical pNJ}
    p_{N,J}^{\mathrm{prac}}(\varepsilon) :=\frac{18V_T^{\mathrm{prac}}}{Nm_{\mathrm{prac}}^2\varepsilon^2} + \frac{\E[e^{\theta B}]}{S_J} \exp\left(-\frac{\theta m_{\mathrm{prac}}}{3}N\varepsilon\right).
\end{equation}
We are now ready to state and prove the practical hit-ratio convergence result.

\begin{reptheorem}{thm: Practical Hit Ratio Convergence Full}[Practical Hit Ratio Convergence]
Fix a finite $J$ such that $S_{J+1}>0$ and $T>T_C^{\mathrm{prac}}$. For every $0<\varepsilon<\min\{T_C^{\mathrm{prac}},T-T_C^{\mathrm{prac}}\},$ the difference between the practical finite-size hit ratio and the mean-field limiting hit ratio can be bounded, such that
\begin{equation*}
    \left|h_{ \mathrm{prac}}^{(N)}-h_{ \mathrm{prac}}^{(\infty)}\right| \leq \frac{L_F\varepsilon A_{H,\mathrm{hash}}^{(J)}}{A_R} + \frac{K_{J,\mathrm{hash}}}{A_R}\sqrt{p_{N,J}^{\mathrm{prac}}(\varepsilon)} + \frac{2\E[B^2e^{\theta B}]e^{-\theta(J+2)}}{A_R}.
\end{equation*}
Consequently,
\begin{equation*}
    h_{ \mathrm{prac}}^{(N)}\rightarrow h_{ \mathrm{prac}}^{(\infty)}= \frac{A_H^{\mathrm{hash}}}{A_R} F(T_C^{\mathrm{prac}}).
\end{equation*}
\end{reptheorem}

\begin{proof}
The same independence argument as in \cref{lem: Exact Representation of LRU Hit Probability} gives
\begin{equation*}
    H_{N,j+1}^{\mathrm{prac}} = \one\left\{W_j<\tau_{C,\mathrm{prac}}^{(N)}(j)\right\},
\end{equation*}
where we recall that the prompt interarrival times $\{W_j\}_{j\geq1}$ are i.i.d. with cumulative distribution function $F$.

Hence we obtain the following practical exact hit ratio representation
\begin{equation}
\label{equ: practical exact hit ratio representation}
    h_{ \mathrm{prac}}^{(N)} = \frac{1}{A_R} \sum_{j\geq1}S_{j+1} \E\left[(R_j-1)F(\tau_{C,\mathrm{prac}}^{(N)}(j))\right].
\end{equation}


Recall that $F$ is $L_F$-Lipschitz-continuous. Repeating the argument in \cref{lem: Finite Turn Hit Ratio Error} with $R_j$ replaced by $R_j-1$, $\tau_C^{(N)}(j)$ replaced by $\tau_{C,\mathrm{prac}}^{(N)}(j)$, and $T_C$ replaced by $T_C^{\mathrm{prac}}$, we obtain that
\begin{align*}
    \left|\sum_{j=1}^{J}S_{j+1}\E\left[(R_j-1)F\left(\tau_{C,\mathrm{prac}}^{(N)}(j)\right)\right]- A_{H,\mathrm{hash}}^{(J)}F(T_C^{\mathrm{prac}}) \right| \leq
    L_F \varepsilon A_{H,\mathrm{hash}}^{(J)}+ K_{J,\mathrm{hash}}\sqrt{p_{N,J}^{\mathrm{prac}}(\varepsilon)}.
\end{align*}

By splitting the infinite series at $J$, and applying the triangle inequality exactly as in the proof of \cref{thm: Hit Ratio Convergence Full}, we obtain
\begin{align*}
    \left|h_{ \mathrm{prac}}^{(N)}-h_{ \mathrm{prac}}^{(\infty)}\right|
    &\leq \frac{L_F \varepsilon A_{H,\mathrm{hash}}^{(J)}}{A_R} + \frac{K_{J,\mathrm{hash}}}{A_R}\sqrt{p_{N,J}^{\mathrm{prac}}(\varepsilon)} + \frac{2\mathcal T_{J,\mathrm{hash}}}{A_R}.
\end{align*}
Applying the tail workload bound \eqref{equ: practical tail workload bound}, we prove \eqref{equ: practical finite N hit ratio bound}.

To prove the hit ratio convergence, we fix $J$ and choose $\varepsilon=N^{-1/4}$, such that $\frac{L_F\varepsilon A_{H,\mathrm{hash}}^{(J)}}{A_R}=O(N^{-1/4})$, and 
\begin{equation*}
    p_{N,J}^{\mathrm{prac}}(N^{-1/4}) = \frac{18V_T^{\mathrm{prac}}}{m_{\mathrm{prac}}^2N^{1/2}} + \frac{\E[e^{\theta B}]}{S_J}\exp\left(-\frac{\theta m_{\mathrm{prac}}}{3}N^{3/4}\right) =O(N^{-1/2}).
\end{equation*}
Hence, for every fixed $J$,
\begin{equation*}
    \limsup_{N\rightarrow\infty} \left|h_{ \mathrm{prac}}^{(N)}-h_{ \mathrm{prac}}^{(\infty)}\right| \leq \frac{2\mathcal T_{J,\mathrm{hash}}}{A_R}.
\end{equation*}
Finally, still by the tail workload bound \eqref{equ: practical tail workload bound}, $\mathcal T_{J,\mathrm{hash}}\rightarrow 0$ as $J\rightarrow\infty$. Taking $J\rightarrow\infty$, we prove \eqref{equ: practical hit ratio convergence full} and therefore \cref{thm: Practical Hit Ratio Convergence Full}.
\end{proof}

\section{Theoretically-Motivated Estimator and LLM Serving Experiments}
\label{sec: LLM-Inference Experiment and Empirical Results}

In this section, we provide a more detailed explanation of the settings and the result analysis, of our empirical LLM serving experiments in \cref{sec: LLM Serving Experiments}. 

Our experiments run Qwen3-8B on five Ascend 910B2 NPUs using a Prefill-Decode disaggregation architecture, with one NPU as the prefiller and four as the decoder pool. Each NPU provides approximately 64 GB of HBM, and the model is deployed under BF16 precision. We reserve 40 GB of the prefiller HBM exclusively for KV cache. With 36 transformer layers, 8 KV heads per layer, and a head dimension of 128, Qwen3-8B requires
\begin{equation*}
    \text{KV Bytes/token}=2\times36\times8\times128\times2=144\text{ KB},
\end{equation*}
where the leading factor $2$ corresponds to the Key (K) and Value (V) matrices. With a block size of 128 tokens, each KV block occupies 18 MB, and an HBM capacity of $N=\lfloor40\text{ GB}/18\text{ MB}\rfloor=2{,}275$ blocks. Our software stack consists of vLLM v0.22.1, vLLM-Ascend v0.22.1rc1, and Mooncake v0.3.9 for KV transfer between the prefill and decode pools.

As mentioned in \cref{sec: LLM Serving Experiments}, we construct the workload from the public ShareGPT dataset (\cite{sharegpt_vicuna_unfiltered}) and sample 6,000 conversations containing at least two turns. The first 1,000 conversations with 4,690 turns are used for cache warm-up and are sufficient to fill the HBM before measurement. The remaining 5,000 conversations form the measurement set. Conversation submissions in both phases follow Poisson processes with some rates specified below. The measurement window starts at the submission time $t_1$ of the first conversation in the measurement set and ends at the submission time $t_{5000}$ of the last. This finite horizon may censor some late turns of active conversations, but we avoid increasing the hit ratio by continuing measurement after new conversation arrivals have stopped.



All block-level workload statistics are computed using the Qwen3 tokenizer and chat template. \cref{fig: E1 Dataset info} shows the empirical conversation-turn distribution and the marginal distribution of new KV blocks per turn. We note that our estimator does not require block increments to be i.i.d. across turns. Using the formulas in equations \eqref{equ: practical estimator AR} to \eqref{equ: practical estimator moments}, we have that 
\[\hat A_R = 38.9092,\quad\hat A_H^{\text{hash}} = 22.5178, \quad \widehat {\E B} = 12.9220,\quad \text{and} \quad\widehat{\E M} = 4.4694.\]


\begin{figure}[ht]
    \centering
    \begin{subfigure}[t]{0.48\linewidth}
        \centering
        \includegraphics[width=\linewidth]{E1_empirical_conversation_turn_distribution.png}
        \caption{Empirical turn number distribution per conversation (Experiment 1)}
        \label{fig: E1_empirical_conversation_turn_distribution}
    \end{subfigure}
    \hfill
    \begin{subfigure}[t]{0.48\linewidth}
        \centering
        \includegraphics[width=\linewidth]{E1_new_kv_block_histogram.png}
        \caption{Marginal distribution of new KV blocks per turn. (Experiment 1)}
        \label{fig: E1_new_kv_block_histogram}
    \end{subfigure}
    \caption{Empirical workload characteristics of the ShareGPT measurement set used in Experiment~1. The measurement set contains 5,000 multi-turn conversations and 22,347 complete user-assistant turns.
    (a): Empirical probability mass function of the number of turns per conversation, with mean 4.4694, variance 7.9659, and maximum 37.
    (b): Marginal empirical probability mass function of the number of new KV blocks per turn. For each turn-$j$, this quantity is calculated by combining the turn-$(j-1)$'s tail tokens and response tokens, and turn-$j$'s whole prompt tokens. The empirical distribution in (b) is used only to describe the workload and does not impose an i.i.d. assumption across turn indices.}
    \label{fig: E1 Dataset info}
\end{figure}

Because ShareGPT provides both prompts and recorded responses, we replay the original conversation text rather than sending generated responses into subsequent turns. For each turn, the decoder is forced to generate the same number of tokens as the corresponding recorded response, preserving the decoding workload and timing, but the generated text and its response KV cache are discarded. This prevents inconsistencies between generated responses and subsequent recorded prompts while preserving the empirical KV-cache workload evolution.

In our experiment, we use the LRU policy equipped in vLLM, which has the following differences from our theoretical model:
\begin{enumerate}
    \item A block's effective recency timestamp is not the arrival time of its latest prompt. Instead, it is determined by system events including completion of prefill-to-decode KV transfer, CPU DRAM write-back, and prefill accesses.
    \item Blocks with nonzero reference counts are temporarily protected from eviction, even if they are currently the LRU blocks.
    \item Blocks belonging to the same conversation do not share identical timestamps. In particular, vLLM exhibits a prefix-preserving, suffix-first eviction behavior, so the cache boundary may partially evict a conversation suffix rather than the entire conversation cohort.
    \item Different conversations may share identical prefix blocks, e.g., through a common system prompt. However, in our experiment, such a cross-conversation prefix sharing is negligible ShareGPT workload used in our experiment.
\end{enumerate}

In Experiment 1, we fix the conversation arrival rate at $\lambda_0=1.5$ and vary the target mean prompt interarrival time $\mu_F$. For turn $j$ of conversation $i$, we let $W_{ij}^{\mathrm{E2E}}$ denote the end-to-end (E2E) time required to complete decoding, and we independently sample $W_{ij}^{\mathrm{sample}}\sim\Exp(1/\mu_F)$. Because the next prompt cannot be submitted before the preceding response is completed, the realized interarrival time is
\begin{equation}
\label{equ: Actual Interarrival time}
    W_{ij}=\max\{W_{ij}^{\mathrm{E2E}},W_{ij}^{\mathrm{sample}}\}.
\end{equation}
The prior estimator uses the target exponential CDF $F(t)=1-\exp(-t/\mu_F)$. After each experiment, we compute the empirical mean interarrival time $\hat\mu_F$ and use the posterior CDF $F^{\mathrm{post}}(t)=1-\exp(-t/\hat\mu_F)$. We report the absolute error $\mathrm{AE}=|\text{estimated hit ratio}-\text{actual hit ratio}|$ and relative error $\mathrm{RE}=\mathrm{AE}/\text{actual hit ratio}\times100\%$.

As shown in \cref{fig: E1-1 hit ratio results}, both estimators correctly capture the decrease in hit ratio as $\mu_F$ increases. 
However, for small $\mu_F$, the event that the interarrival time is not exponentially distributed (i.e., $W^{\text{E2E}}_{ij} > W^{\text{sample}}_{ij}$) occurs more frequently, which makes the empirical mean interarrival time $\hat \mu_F > \mu_F$. In contrast, for large $\mu_F$, the finite measurement horizon censors conversations with long interarrival times, which can make $\hat\mu_F<\mu_F$. These effects explain the relative ordering of the prior and posterior estimates across different $\mu_F$ regimes. Nevertheless, \cref{fig: E1-1 Error Analysis} shows that the absolute error remains below $0.02$ throughout the experiment, while the relative error remains below $16\%$. When the target and empirical interarrival means are close, i.e., $120\leq\mu_F\leq165$ s, the relative errors of both estimators are below $10\%$.



\begin{figure}[ht]
    \centering
    \includegraphics[width=0.6\linewidth]{E1_actual_prior_posterior_hit_ratio.png}
    \label{fig: E1_actual_prior_posterior_hit_ratio}
    \caption{The plot of the actual and the prior/posterior estimated hit ratios as functions of target mean interarrival times $\mu_F$ (Experiment 1). We use 5000 ShareGPT real prompt-response conversations for measurement and 1000 for warm-up. The arrival process is a Poisson process with rate $1.5$ per second, and the interarrival time is sampled according to \eqref{equ: Actual Interarrival time}, with target means chosen from $90$s to $240$s. The measurement duration is between the submission times of the 1-st and the 5000-th measurement conversations.
    }
    \label{fig: E1-1 hit ratio results}
\end{figure}

\begin{figure}[ht]
    \centering
    \begin{subfigure}[t]{0.48\linewidth}
        \centering
        \includegraphics[width=\linewidth]{E1_posterior_estimator_errors_AE.png}
        \caption{Absolute errors. (Experiment 1)}
        \label{fig: E1_posterior_estimator_errors_AE}
    \end{subfigure}
    \hfill
    \begin{subfigure}[t]{0.48\linewidth}
        \centering
        \includegraphics[width=\linewidth]{E1_posterior_estimator_errors_RE.png}
        \caption{Relative errors. (Experiment 1)}
        \label{fig: E1_posterior_estimator_errors_RE}
    \end{subfigure}
    \caption{The absolute and relative errors of the prior/posterior estimated hit ratios as functions of $\mu_F$ (Experiment 1).}
    \label{fig: E1-1 Error Analysis}
\end{figure}


In Experiment 2, to reduce the error caused by the mismatch between the target and realized interarrival distributions, we fix $\mu_F=135$ s, for which $\hat\mu_F$ is closest to the target in Experiment 1. We vary the conversation arrival rate from $0.5$ to $2.0$. As shown in \cref{fig: E1-2 hit ratio results}, both estimators correctly capture the decrease in empirical hit ratio as $\lambda_0$ increases. Moreover, the estimates are more accurate in this experiment. To be specific, we note in \cref{fig: E1-2 Error Analysis} that both estimators achieve absolute errors below $0.015$ and relative errors below $7\%$.

\begin{figure}[ht]
    \centering
    \includegraphics[width=0.6\linewidth]{E1-2_arrival_rate_actual_prior_posterior_hit_ratio.png}
    \label{fig: E1-2_arrival_rate_actual_prior_posterior_hit_ratio}
    \caption{The plot of the actual and the prior/posterior estimated hit ratios as functions of the conversation arrival rate $\lambda_0$ (Experiment 2). We use 5000 ShareGPT real prompt-response conversations for measurement and 1000 for warm-up. The arrival process is a Poisson process with rate chosen from 0.5 to 2.0, and the interarrival time is sampled according to \eqref{equ: Actual Interarrival time}, with target mean $135$s.
    }
    \label{fig: E1-2 hit ratio results}
\end{figure}

\begin{figure}[ht]
    \centering
    \begin{subfigure}[t]{0.48\linewidth}
        \centering
        \includegraphics[width=\linewidth]{E1-2_arrival_rate_estimator_errors_AE.png}
        \caption{Absolute errors. (Experiment 2)}
        \label{fig: E1-2_arrival_rate_estimator_errors_AE}
    \end{subfigure}
    \hfill
    \begin{subfigure}[t]{0.48\linewidth}
        \centering
        \includegraphics[width=\linewidth]{E1-2_posterior_estimator_errors_RE.png}
        \caption{Relative errors. (Experiment 2)}
        \label{fig: E1-2_posterior_estimator_errors_RE}
    \end{subfigure}
    \caption{The absolute and relative errors of the prior/posterior estimated hit ratios as functions of $\lambda_0$ (Experiment 2).
    }
    \label{fig: E1-2 Error Analysis}
\end{figure}

\end{document}
